%%% POST-RECORD v1 BEGIN
% format: "hexagon-post-format/1"
% record_id: "Q03"
% version: 1
% date: "2026-11-15"
% title: "Polynomial-time unit-job scheduling for every fixed number of machines m >= 3"
% kind: "generalization"
% claim_status: "proved"
% statement_plain: "For every fixed integer m >= 3, explicitly listed unit jobs with arbitrary acyclic precedence on m identical machines admit exact deadline decision, witness construction and optimum construction in at most G(m)(L+2)^(60m^2+129m+2) deterministic fixed-multitape steps. A parsed canonical padded deadline with N=mT jobs uses at most A(m,12m+25)N^(60m^2+129m+1) steps. The description hierarchy has 12m+25 global leaves and 6m+16 relative-prefix leaves."
% statement_tex: "For each fixed $m\\geq3$, $Pm\\mid\\mathrm{prec},p_j=1\\mid C_{\\max}$ is solvable in deterministic multitape time $O_m((L+2)^{60m^2+129m+2})$; a parsed padded deadline uses $O_m(N^{60m^2+129m+1})$ steps, $N=mT$."
% extends: "OpenAI, preprint of 24 September 2026, Theorem 1.1 (three machines); all steps are proved here."
% known_before: "Two machines: Fujii-Kasami-Ninomiya with erratum, Coffman-Graham; variable-machine NP-completeness: Ullman; three machines: OpenAI; four-machine theorem code: drbh bbd3013; all-fixed-m theorem and linear budgets stated by Jansen arXiv:2611.09392v1."
% new_here: "Self-contained audited capacity-m proof and itemized costs of the supplied Q03 construction; no novelty or priority claim."
% main_statements: ["Theorem A = Theorem 1.1 of the supplied return, established in 9.1-9.3", "Theorem B = Theorem 8.5"]
% statements_in_proof: 44
% hard_step: "Lemma 5.5; Lemmas 6.2 and 6.5; Lemma 7.1; Lemma 8.4"
% external_results: "none load-bearing"
% produced_by: [{"role": "question", "model": "Claude Opus 5.5", "effort": "effort not recorded", "date": "2026-10-22"}, {"role": "proof return", "model": "GPT-6 Astra", "effort": "ChatGPT Pro mode", "date": "2026-10-23"}, {"role": "format text", "model": "Claude Fable 5.2", "effort": "high", "date": "2026-11-14"}, {"role": "conversion", "tool": "pandoc 3.9.0.2", "date": "2026-11-14"}, {"role": "deposit-task author", "model": "Claude Opus 5.5", "effort": "max", "date": "2026-11-15"}, {"role": "current audit and editing", "model": "GPT-6 Astra", "effort": "ChatGPT Pro mode", "date": "2026-11-15"}]
% checks: [{"id": "H1", "by": "Claude Opus 5.5", "effort": "max, Message Batches API", "date": "2026-10-23", "type": "F,D", "scope": "supplied return", "verdict": "proved with fixable gaps"}, {"id": "H2", "by": "Claude Opus 5.5", "effort": "max, Message Batches API", "date": "2026-10-23", "type": "D", "scope": "supplied return", "verdict": "proved as stated"}, {"id": "H3", "by": "Claude Opus 5.5", "effort": "max, Message Batches API", "date": "2026-10-23", "type": "F", "scope": "supplied return", "verdict": "proved as stated"}]
% checks_of_this_text: [{"by": "GPT-6 Astra (OpenAI)", "effort": "ChatGPT Pro mode", "date": "2026-11-15", "type": "F,D,C,S", "scope": "all 44 numbered statements in Sections 2-9; main theorems; rewritten proofs; costs; source attribution; finite checks; post consistency", "verdict": "proved as stated", "limits": "S-FKN original full text unavailable; S-L archival endpoint 404; historical prompt manifest discrepancy; no independent Lean build"}]
% human_understanding: "none"
% text_origin: "primarily AI-generated"
% formalization: "none for this proof"
% open_points: ["No mathematical gap found in Sections 2-9. The three original cost/source audit items are addressed by the recorded checks.", "Original FKN 1969 full text unavailable; publisher metadata/abstract and 1971 erratum opened.", "Course provenance.json endpoint returned HTTP 404; first-publication dates are unverified.", "Supplied original prompt is 9551 bytes; historical manifest lists an absent 9383-byte variant with a different hash. Supplied prompt preserved verbatim.", "Historical run dates, API mode and independence are supplied-record metadata; raw API logs absent.", "OpenAI/drbh Lean builds and Jansen/Giraudin full proofs were not independently replayed or audited; their declared scope was checked."]
% not_claimed: ["a bound polynomial in m", "a practical algorithm", "release dates, individual deadlines, unequal machine speeds or unequal job lengths", "formalization of this proof", "novelty or priority", "optimality of the displayed exponent"]
% literature_search: {"date": "2026-11-15", "scope": "arXiv; STOC, FOCS, SODA, ICALP, IPCO, ESA, APPROX; authors pages; GitHub; general web; Hexagon; title/citation searches for the OpenAI preprint", "outcome": "found drbh four-machine code, Jansen all-fixed-m theorem statement and linear budgets, Giraudin three-machine exponent statement; no fixed-m>=4 hardness result located; no priority claim"}
% licence: "CC0-1.0"
% provenance: {"prompt": "anc/ai-provenance/Q03_PROMPT.txt", "current_request": "anc/Q03_DEPOSIT_prompt.md", "return": "anc/ai-provenance/Q03-solution.md", "return_sha256": "3ebeb452d6f461c23b4f8fa6f04ba6ed15abac541b5d2562fa814ba77977e60c", "reports": ["anc/ai-provenance/Q03-check.md", "anc/ai-provenance/Q03-b-check-b.md", "anc/ai-provenance/Q03-c-check-c.md"], "current_report": "Q03_DEPOSIT_report.md", "repository": "ancillary files supplied with this deposit package; no external publication by this audit"}
% hexagon: {"subject": "cs.DS", "cross_lists": ["cs.DM", "cs.CC", "math.CO", "math.OC"], "msc": ["90B35", "68Q25", "68W40", "06A07"], "tags": ["scheduling", "precedence constraints", "unit jobs", "identical machines", "makespan", "fixed number of machines", "polynomial-time algorithm", "dynamic programming", "partially ordered sets"], "tools": ["GPT-6 Astra, ChatGPT Pro mode", "Claude Opus 5.5, effort max", "Claude Opus 5.5, effort max, Message Batches API", "Claude Opus 5.5, effort not recorded", "Claude Fable 5.2, effort high", "pandoc 3.9.0.2", "Python 3", "pdfLaTeX"], "comments": "Current text audited F,D,C,S by GPT-6 Astra, ChatGPT Pro mode; cost exponent corrected to 60m^2+129m+2. Source-access and archival qualifications stated. No human proof check or formalization of this text. No novelty/priority claim.", "provenance_statement": "See the Provenance and Declarations parts of the post, reproduced in the accompanying form answers."}
%%% POST-RECORD v1 END
\documentclass[11pt,a4paper]{article}
\usepackage[T1]{fontenc}
\usepackage{lmodern,microtype}
\usepackage[margin=2.4cm]{geometry}
\usepackage{amsmath,amssymb,amsthm,mathrsfs}
\usepackage{booktabs,tabularx,longtable,array,calc}
\usepackage{enumitem,xcolor,needspace}
\usepackage{titling}
\setlength{\droptitle}{-1.5\baselineskip}
\PassOptionsToPackage{hyphens}{url}
\usepackage[hidelinks]{hyperref}
\emergencystretch=2em
\providecommand{\tightlist}{\setlength{\itemsep}{0pt}\setlength{\parskip}{0pt}}
\theoremstyle{plain}
\newtheorem{maintheorem}{Theorem}
\renewcommand{\themaintheorem}{\Alph{maintheorem}}
\theoremstyle{definition}
\newtheorem{cited}{Cited result}
\newenvironment{numbered}[3]{\par\addvspace{\medskipamount}\Needspace{3\baselineskip}\noindent
  \textbf{#1 #2}\if\relax\detokenize{#3}\relax\else\ (#3)\fi\textbf{.}\ \itshape\ignorespaces}
  {\par\addvspace{\smallskipamount}}
\newenvironment{numberedup}[3]{\par\addvspace{\medskipamount}\Needspace{3\baselineskip}\noindent
  \textbf{#1 #2}\if\relax\detokenize{#3}\relax\else\ (#3)\fi\textbf{.}\ \ignorespaces}
  {\par\addvspace{\smallskipamount}}
\newcommand{\uses}[1]{\par\nopagebreak[4]\noindent{\small\textsf{Uses:} #1}\par\smallskip}
\newcommand{\step}[2]{\par\smallskip\noindent\textbf{Step #1} (#2).\ \ignorespaces}
\newcommand{\by}[1]{{\small[#1]}}
\newcolumntype{L}[1]{>{\raggedright\arraybackslash}p{#1}}
\newcolumntype{Y}{>{\raggedright\arraybackslash}X}
\newcommand{\fixedhead}[1]{\section*{#1}}
\title{Polynomial-time unit-job scheduling\\for every fixed number of machines $m\geq3$}
\author{Human contributor: entered at deposit}
\date{15 November 2026 \quad Version 1 \quad CC0 1.0}
\begin{document}
\maketitle
\vspace{-1.5\baselineskip}

\begin{abstract}
For every fixed $m\geq3$, unit-job scheduling with explicitly listed precedence constraints on $m$ identical machines admits exact deadline decision and optimal-schedule construction in $O_m((L+2)^{60m^2+129m+2})$ deterministic multitape steps; a parsed, padded deadline of size $N=mT$ costs $O_m(N^{60m^2+129m+1})$.
OpenAI proved the three-machine result (2026); drbh supplied four-machine Lean code (2026); Jansen states the fixed-machine result in arXiv:2611.09392 (2026).
This generalization gives a self-contained proof with leaf budget $12m+25$ and an itemized tape-cost derivation, without a priority claim.
A finite dynamic program uses Boolean descriptions; an injection into one boundary slot bounds the lists, and local restriction identities control direction changes.
GPT-6 Astra, ChatGPT Pro mode, produced the proof and the current step-by-step, computational and source audit. Three earlier reports are by Claude Opus 5.5, effort max, Message Batches API. Human understanding: none.
Source-access and archival qualifications are recorded in Limits and open points.
\end{abstract}

\section{Result}
\subsection{Statement}

The input explicitly lists a DAG $G=(V,E)$, $n=|V|\geq1$ jobs, and binary integers $m\geq3,T$.
A schedule is a map $\tau:V\to\{1,\ldots,T\}$ with at most $m$ jobs per slot and $\tau(u)<\tau(v)$ on every arc.
The makespan is $\max\tau$. The canonical encoding and its bit length $L$ are specified in Section~9.

\begin{cited}[OpenAI, {\cite[Theorem~1.1]{P}}]\label{cit:base}
For $m=3$, deadline decision and witness construction for $1\leq T\leq n$, and optimum-schedule construction, take polynomial deterministic multitape time in the explicit input length.
\end{cited}
\uses{Source statement opened and compared in source check S-P; the proof below is self-contained.}

\noindent\begin{tabularx}{\linewidth}{@{}lYll@{}}
\toprule
ID & Result & Status & Proof\\
\midrule
A & Exact algorithm for every fixed $m\geq3$ & proved & 9.1--9.3\\
B & Hierarchy with $12m+25$ global leaves & proved & 8.5\\
\bottomrule
\end{tabularx}

\enlargethispage{\baselineskip}
\begin{maintheorem}[Theorem 1.1 of the checked return]\label{thm:A}
One deterministic algorithm, given the input above, decides deadline feasibility, constructs a feasible schedule on acceptance, and constructs an optimal schedule. For each fixed $m$ all these tasks use
\[
 O_m\bigl((L+2)^{60m^2+129m+2}\bigr)\tag{1.1}
\]
steps on a deterministic Turing machine with a fixed number of tapes.
For $1\leq T<n$ and $n\leq mT$, its deadline computation on the parsed, canonically numbered padded DAG with $N=mT$ jobs, including its witness, uses
\[
 O_m\bigl(N^{60m^2+129m+1}\bigr).\tag{1.2}
\]
Proposition 9.2 and Corollary 9.3 give explicit effective coefficient bounds $A(m,12m+25)$ and $G(m)$, respectively.
\end{maintheorem}
\uses{Theorem 9.1; Proposition 9.2; Corollary 9.3.}
\pagebreak

\begin{maintheorem}[Theorem 8.5 of the checked return]\label{thm:B}
For $m\geq3$, every full schedule of the padded job set $J$, $|J|=mT$, admits a description hierarchy as in Definition 2.9. The root has its one-leaf description $J$, every other node has a global description of at most $12m+25$ leaves, and every marked original-time prefix has a relative description of at most $6m+16$ leaves, independently of depth.
\end{maintheorem}
\uses{Definition 2.9; Proposition 7.3; Lemmas 8.2--8.4.}

\noindent\textbf{Instance.} For $m=4$, $K_4=73$, (1.2) has exponent $1477$, and (1.1) has exponent $1478$.
\textbf{Not claimed:} a bound polynomial in $m$; a practical algorithm; extensions to release dates, individual deadlines, unequal machine speeds or unequal job lengths; formalization of this proof.

\subsection{Status}
\noindent\begin{tabularx}{\linewidth}{@{}L{3.0cm}Y@{}}
\toprule
Kind & generalization\\
Claim & proved\\
Checks & Current text: F, D, C, S by GPT-6 Astra, ChatGPT Pro mode, 15 November 2026; verdict: proved as stated. Supplied return: three earlier reads by Claude Opus 5.5, effort max, Message Batches API.\\
Human understanding & none\\
Formalization & none for this proof\\
Open points & No mathematical gap found in Sections 2--9; source-access and archival qualifications are listed below.\\
\bottomrule
\end{tabularx}

\subsection{Known and new}
\noindent\begin{tabularx}{\linewidth}{@{}L{4.0cm}Y@{}}
\toprule
Statement & Opened source and scope\\
\midrule
Two machines & Coffman--Graham \cite[Theorem~1, p.~203]{CG}; Fujii--Kasami--Ninomiya \cite[abstract]{FKN} with the correction \cite[p.~141]{FKNe}.\\[2pt]
Machine count in the input & Ullman's unit-time NP-completeness result \cite[Theorem~1, p.~389]{U}; the fixed-count question is discussed on p.~392.\\[2pt]
Three machines & OpenAI \cite[Theorem~1.1]{P}. The earlier question is OPEN8 \cite[Appendix A13, p.~287]{GJ}. Giraudin's subsequent statement has input exponent $276$ \cite[Theorem~1]{G}.\\[2pt]
Four machines & drbh's commit \texttt{bbd3013}, \cite[\texttt{Model.lean}, \texttt{ReleaseTheorem}; \texttt{Main.lean}, \texttt{main\_theorem}]{D}, contains the capacity-four solver and theorem code.\\[2pt]
Every fixed $m\geq3$ & Jansen \cite[Theorem~1.1]{J} states exact decision and optimization; \cite[Lemma~7.15]{J} gives node/cut leaf budgets $4m+5$ and $m+3$.\\[2pt]
Exact subexponential bound & Nederlof--Swennenhuis--W\k{e}grzycki \cite[Theorem~1.1]{N}: $(1+n/m)^{O(\sqrt{nm})}$.\\[2pt]
Approximation scheme & Das--Wiese \cite[Theorem~1]{A}: $(1+\varepsilon)$ approximation with the parameter-dependent quasipolynomial bound stated there.\\[2pt]
This post & Theorems A and B are proved here with $12m+25$ leaves and the derived exponent in (1.1). No novelty or priority is claimed.\\
\bottomrule
\end{tabularx}

\smallskip
\noindent\textbf{Relation of the guarantees.}
Lemma 2.2 gives integer makespans, and a topological schedule gives $\mathrm{OPT}\leq n$.
Using $\varepsilon=1/(n+1)$ in the guarantee of \cite[Theorem~1]{A} yields a feasible integer makespan strictly below $\mathrm{OPT}+1$, hence exactly $\mathrm{OPT}$.
The source's time bound depends on this varying $\varepsilon$; (1.1) is derived directly from the finite search in this post.
For fixed $m$, the displayed bound of \cite[Theorem~1.1]{N} has the form $\exp(O_m(\sqrt n\log(n+1)))$, by taking its logarithm. This is the stated upper bound, with no claim of a corresponding lower bound on that algorithm.

\smallskip
\noindent\textbf{Credit and search.}
The construction follows OpenAI \cite[Sections 2--9]{P} and the course \cite[lessons 1--8]{L}.
The exact scope of drbh's four-machine code is credited in the source record; its Lean build was not replayed in this audit.
The search of 15 November 2026 covered arXiv, STOC, FOCS, SODA, ICALP, IPCO, ESA, APPROX, authors' pages, GitHub, the general web and Hexagon, including citations and title searches for [P]. It found [D], [J] and [G]; no fixed-$m\geq4$ hardness result was located. Every query and outcome is in \texttt{Q03\_DEPOSIT\_literature.md}.


\subsection{Proof map}

\noindent\textbf{Idea.} Pad to full $m$-job slots and enumerate bounded Boolean descriptions of job sets. A state path joins smaller accepted sets through single slots. Separators split a feasible schedule into a hierarchy. The map of Lemma 5.5 sends surviving injection entries injectively into one boundary slot; the restricted equalities of Lemma 6.1 keep their descriptions bounded through direction changes.

\begin{enumerate}[nosep]
\item Foundations, finite search and its certificate: Lemmas 2.1--2.7, Proposition 2.8, Definition 2.9 and Lemma 2.10.
\item Exchanges, separators and walks: Lemmas 3.1--3.6 and Corollary 3.7.
\item Centers, sides and child labels: Lemmas 4.1--4.5.
\item Carried lists and the capacity injection: Lemmas 5.1--5.5.
\item Direction changes and formula counts: Lemmas 6.1--6.5.
\item Exterior invariants and induction: Lemmas 7.1--7.2 and Proposition 7.3.
\item Global child identities, prefixes and the budget: Lemmas 8.1--8.4 and Theorem 8.5.
\item Correctness and itemized multitape costs: Theorem 9.1, Proposition 9.2 and Corollary 9.3; compact decision input: Proposition 9.4.
\end{enumerate}

\noindent\textbf{Hard steps.} Lemma 5.5 proves the $m$-entry injection bound; Lemmas 6.2 and 6.5 prevent leaf counts from growing with depth; Lemma 7.1 propagates exterior qualifiers through the first advance using (I3). Their independent derivations are in the current audit. The constants are
\[
 A_m=3m+(3m+7)+2=6m+9,\qquad
 K_m=(A_m+2)+(A_m+5)=12m+25.
\]

\noindent\textbf{Reading routes.} The result and model are in Section 1.1 and Section 9. Readers comparing with the three-machine proof can follow 3.3, 5.5, 6.2, 6.5, 7.1, 8.4 and 9.2. The full proof runs through Sections 2--9 in order; Proposition 7.3 checks the construction's induction order.

\noindent\textbf{External results used.} None load-bearing. Source checks compare the construction with OpenAI's preprint and the course; every needed mathematical step is proved here. The external results in the preceding table provide historical and algorithmic context.

\noindent\textbf{Dependencies.} The following table covers all 44 numbered statements of Sections 2--9, including Definition 2.9. The source column is checked against the pinned public source of [P]. Its entries identify the three-machine counterparts; the last column records the change proved here. Definitions and displayed constructions supplement the numbered dependencies where the Uses line says so.

{\footnotesize
\begin{longtable}{@{}L{2.6cm}L{3.2cm}L{3.7cm}L{5.1cm}@{}}
\caption{All numbered proof statements, source counterparts and direct dependencies.}\label{tab:dep}\\
\toprule Statement & Source counterpart [P], S-P & Uses & Capacity transfer\\\midrule\endfirsthead
\toprule Statement & Source counterpart [P], S-P & Uses & Capacity transfer\\\midrule\endhead
\bottomrule\endfoot
Lemma 2.1 & Section 2.2 (prose) & definitions only. & Unchanged argument: Transitivity and strict inequalities only. Finite acyclicity only.\\[3pt]
Lemma 2.2 & Section 1 (prose) & Lemma 2.1; definitions of a unit-job schedule and machine assignment. & Changed: Machine assignment has target \([m]\); augmented arcs give integer starts.\\[3pt]
Lemma 2.3 & Lemma 2.1; \textsf{lem:\allowbreak padding} & definitions only. & Changed: Exactly \(mT-n\) isolated jobs and \(N\leq mn\).\\[3pt]
Lemma 2.4 & Lemma 2.2; \textsf{lem:\allowbreak reorder} & definitions only. & Unchanged argument: Equal interval length \(|W|/m\); exterior relations unchanged.\\[3pt]
Lemma 2.5 & Section 2.2 (prose) & definitions only. & Unchanged argument: Antichain property, not the number of companions.\\[3pt]
Lemma 2.6 & Section 2.3 (prose) & definitions only. & Unchanged argument: De Morgan transformation preserves leaves.\\[3pt]
Lemma 2.7 & Lemma 2.3; \textsf{lem:\allowbreak family} & Lemma 2.1; definitions of the atomic alphabet and formula evaluation. & Changed: \(\binom Nm\leq N^m\) replaces \(\binom N3\).\\[3pt]
Proposition 2.8 & Proposition 3.1; \textsf{prop:\allowbreak algorithm} & definitions of the state graph, compatibility, and acceptance in Section 2.2; induction on set cardinality. & Changed: One nonempty \(m\)-slot per state; compatibility checks all relations.\\[3pt]
Definition 2.9 & Definition 3.2; \textsf{def:\allowbreak hierarchy} & definitions of full schedules, open intervals, and global and relative descriptions. & Changed: Slot size $3$ becomes $m$; the hierarchy clauses are retained.\\[3pt]
Lemma 2.10 & Lemma 3.3; \textsf{lem:\allowbreak hierarchy-\allowbreak dp} & Definition 2.9; state and acceptance definitions in Section 2.2. & Unchanged argument: Actual marked slots form an accepting state path.\\[3pt]
Lemma 3.1 & Lemma 4.2; \textsf{lem:\allowbreak separator-\allowbreak description} & Lemma 2.5; separator definition and (3.1). & Unchanged argument: The two separator conditions and weak-cone exclusion.\\[3pt]
Lemma 3.2 & Lemma 4.2; \textsf{lem:\allowbreak separator-\allowbreak description} & definitions of frame reversal, cutoff, and separator. & Unchanged argument: Complement cutoff, reversed time and precedence.\\[3pt]
Lemma 3.3 & Lemma 4.3; \textsf{lem:\allowbreak lex-\allowbreak prefix} & the definition of a normalized schedule (Section 3.1). & Changed: Every sorted order statistic weakly decreases; one decreases strictly.\\[3pt]
Lemma 3.4 & Lemma 4.4; \textsf{lem:\allowbreak shared} & Lemma 3.3; definitions of domain, priority prefix, exception, and separator. & Unchanged argument: Swapping two jobs and tracing increasing successor times.\\[3pt]
Lemma 3.5 & Lemma 4.4; \textsf{lem:\allowbreak shared} & Lemma 3.4. & Unchanged argument: Strictly increasing finite sequence of positions.\\[3pt]
Lemma 3.6 & Lemma 4.4; \textsf{lem:\allowbreak shared} & Lemmas 3.4 and 3.5. & Unchanged argument: Only the promoted job changes status; subsequent exceptions are old low jobs.\\[3pt]
Corollary 3.7 & Corollary 4.5; \textsf{cor:\allowbreak any-\allowbreak ideal} & Lemma 3.6; definitions of cutoff and normalization. & Unchanged argument: Priority-prefix promotions from a normalized schedule.\\[3pt]
Lemma 4.1 & Section 5.1 (prose) & definitions only. & Unchanged argument: Global upsets and increasing rank.\\[3pt]
Lemma 4.2 & Section 5.2 (prose) & invariant (I1); frame and cone definitions. & Unchanged argument: (I1), complements and reversal.\\[3pt]
Lemma 4.3 & Section 5.2 (prose) & Lemmas 4.1 and 4.2; invariant (I1). & Unchanged argument: Downset blocks and key order.\\[3pt]
Lemma 4.4 & Lemma 5.1; \textsf{lem:\allowbreak edge-\allowbreak labels} & Lemmas 3.6 and 4.3; the center and side-walk construction. & Unchanged argument: Shared walk plus center separator.\\[3pt]
Lemma 4.5 & Lemma 5.1; \textsf{lem:\allowbreak edge-\allowbreak labels} & Lemmas 4.3 and 4.4; Definition 2.9 and the child construction. & Unchanged argument: Genuine right boundary, low/high cone inclusions, strict shrinking.\\[3pt]
Lemma 5.1 & Lemma 6.1; \textsf{lem:\allowbreak carry} & definitions only. & Unchanged argument: Three cases for the largest qualifying index, unrelated to capacity.\\[3pt]
Lemma 5.2 & Lemma 6.2; \textsf{lem:\allowbreak prune} & definitions only. & Unchanged argument: The smaller-key job must belong to the pruning domain.\\[3pt]
Lemma 5.3 & Lemma 6.3; \textsf{lem:\allowbreak update-\allowbreak labels} & Lemmas 4.2, 4.3, 4.5, and 5.1; invariant (I1); the stipulated switch-upset property (5.3). & Unchanged argument: Last injection index and ordered switch pair.\\[3pt]
Lemma 5.4 & Section 6.3 (prose) & definitions of carry, pruning, and a run. & Unchanged argument: Suffix thresholds combine by their maximum.\\[3pt]
Lemma 5.5 & Lemma 6.4; \textsf{lem:\allowbreak list-\allowbreak length} & Lemmas 4.5, 5.1, 5.3, and 5.4; invariant (I1); injection definition (5.2). & Changed: Explicit injective function into a boundary with \(m\) jobs.\\[3pt]
Lemma 6.1 & Lemma 6.5; \textsf{lem:\allowbreak flatten} & Lemmas 5.1 and 5.4; the pruning and run definitions; switch pair (5.4). & Unchanged argument: While the first entry survives, the second is untrimmed or locally absent.\\[3pt]
Lemma 6.2 & Lemma 6.6; \textsf{lem:\allowbreak old-\allowbreak low} & Lemmas 2.6, 5.3, 5.4, 5.5, and 6.1; the edge-cutoff construction. & Changed: \(3m+3\) leaves; only equality on the earlier node is needed.\\[3pt]
Lemma 6.3 & Section 6.5 (prose) & Lemma 4.5; the frame-reversal map. & Unchanged argument: Reverse the old side inequalities.\\[3pt]
Lemma 6.4 & Lemma 6.7; \textsf{lem:\allowbreak G} & Lemmas 2.5, 4.3, 4.4, 4.5, 6.2, and 6.3; definition (6.8). & Unchanged argument: Four upset factors; localize to the old node before using \(H^*\).\\[3pt]
Lemma 6.5 & Lemma 6.8; \textsf{lem:\allowbreak entry-\allowbreak size} & Lemmas 2.6, 5.4, 5.5, and 6.2; definitions (5.2), (5.4), and (6.8). & Changed: At most \(3m+7\) leaves per entry; aggregate \(6m+9\).\\[3pt]
Lemma 7.1 & Lemma 7.1; \textsf{lem:\allowbreak past} & Lemmas 2.5, 3.4, 3.6, 4.3, 4.4, 5.1, 5.3, and 6.4; invariants (I1), (I2), and (I3). & Unchanged argument: First advance uses parent (I3); later advances use Lemma 3.4.\\[3pt]
Lemma 7.2 & Lemma 7.2; \textsf{lem:\allowbreak swap} & Lemmas 2.4, 2.5, 3.3, 5.1, 5.2, 5.3, and 6.4; invariants (I1) and (I3). & Unchanged argument: A child reordering changes only a suffix of normalized \(E\).\\[3pt]
Proposition 7.3 & Proposition 7.3; \textsf{prop:\allowbreak construction-\allowbreak valid} & Lemma 2.4; Corollary 3.7; Lemmas 4.3--4.5, 5.3, 5.5, 6.1--6.5, 7.1, and 7.2; the root invariants. & Unchanged argument: Induction uses only node/ancestor data; no descendant assumption.\\[3pt]
Lemma 8.1 & Section 8.1 (prose) & Lemmas 2.6, 4.3, 5.1, and 6.5; definitions (8.1) and (8.2). & Changed: The same restriction identities use the capacity-dependent carry bound $A_m=6m+9$.\\[3pt]
Lemma 8.2 & Lemma 8.1; \textsf{lem:\allowbreak global-\allowbreak child} & Lemmas 2.5, 4.3, 4.4, 4.5, 5.3, 7.1, and 8.1; the center-separator property. & Unchanged argument: Exact global equalities (8.4)--(8.5), with no parent formula.\\[3pt]
Lemma 8.3 & Lemma 8.2; \textsf{lem:\allowbreak relative-\allowbreak splits} & Lemmas 2.6, 3.1, 4.4, and 8.1; the center-cutoff definition. & Unchanged argument: Relative separator descriptions and subtraction of the reversed slot.\\[3pt]
Lemma 8.4 & Lemma 8.3; \textsf{lem:\allowbreak final-\allowbreak count} & Lemmas 2.6, 6.5, 8.1, 8.2, and 8.3; the atomic formula alphabet. & Changed: Linear aggregate accounting: \(12m+25\) and \(6m+16\).\\[3pt]
Theorem 8.5 & Theorem 3.4; \textsf{thm:\allowbreak structure} & Definition 2.9; Proposition 7.3; Lemmas 8.2, 8.3, and 8.4. & Changed: Global budget \(K_m=12m+25\).\\[3pt]
Theorem 9.1 & Theorem 1.1; \textsf{thm:\allowbreak main} & Lemmas 2.1, 2.2, 2.3 and 2.10; Proposition 2.8; Theorem 8.5. & Changed: Proved for every fixed \(m\geq3\), with an explicit exponent.\\[3pt]
Proposition 9.2 & Section 9 (prose) & Lemmas 2.1 and 2.7; Proposition 2.8; the state and transition definitions in Section 2. & Changed: $B=K512^KN^{mK}$; $\mathcal{S}=BN^m$; every tape subroutine is counted.\\[3pt]
Corollary 9.3 & Section 9 (prose) & Lemma 2.3; Lemma 8.4 and Theorem 8.5; Theorem 9.1; Proposition 9.2. & Changed: $N\leq mL$ and at most $L$ trials give exponent $60m^2+129m+2$.\\[3pt]
Proposition 9.4 & Remark 9.1; \textsf{rem:\allowbreak compact} & Lemmas 2.1 and 2.3; Theorem 9.1; Proposition 9.2. & Changed: Explicit compact expansion $n\leq N<2m|E|$; decision output is bounded.\\[3pt]
\end{longtable}}

\subsection*{Notation}

\noindent
{\footnotesize
\begin{longtable}{@{}L{4.0cm}L{8.2cm}L{2.5cm}@{}}
\toprule
Symbol or term & Meaning & Defined in \\
\midrule
\endhead
\bottomrule
\endfoot
$m$; $G=(V,E)$, $n$, $T$, $L$ & number of machines, fixed, $m\geq3$; the instance, its number of jobs, the deadline, the binary input length & Section 1.1 \\
$x\prec y$ & there is a directed path of positive length from $x$ to $y$ & Lemma 2.1 \\
$J$, $N=mT$ & the padded job set and its size & Lemma 2.3 \\
$\lambda\colon J\to[N]$ & rank: the position in the deterministic total order extending $\prec$ & after Lemma 2.1 \\
full schedule of $W$ & a sequence of $m$-element antichains partitioning $W$ and respecting $\prec$ & Lemma 2.3 \\
frame (forward, reverse) & the direction in which precedence, time and rank are read & Section 2.1 \\
$P_a$, $D_a$, $\widehat P_a$, $\widehat D_a$ & strict predecessor and successor cones of an $m$-element set $a$; with a hat, united with $a$ & Section 2.1 \\
upset, downset & subsets closed under successors, under predecessors & Section 2.1 \\
formula, leaf, $[F]$, $\mathcal F_K$ & a Boolean tree over the atomic tests; its evaluated set; the sets with a formula of at most $K$ leaves & Section 2.1 \\
global, relative description & $W\in\mathcal F_K$; $L=W\cap[F]$ & Section 2.1 \\
state $(L,Z,R)$; $\mathsf A$ & a state of the dynamic program; its acceptance function & Section 2.2 \\
hierarchy, node, mark & the certificate of completeness & Definition 2.9 \\
domain $U$; cutoff $S$; low, high & job set of an open interval; a relative downset of $U$; in $S$, outside $S$ & Section 3 \\
separator at $S$; (a), (b) & a slot $z$ with every earlier high job in $P_z$ and every later low job in $D_z$ & Section 3 \\
$B_z(H)$ & $P_z\cup(H\setminus\widehat D_z)$ & (3.1) \\
priority order $\pi$; normalized & a total order extending $\prec$; lexicographically minimal encoded schedule & Section 3.1 \\
exception for $L$ at $s$ & a job of $L$ after $s$ outside $D_s$ & Lemma 3.4 \\
$W$; $a$, $b$; tag & a node; its boundaries; its frame & Section 4 \\
$\mathcal K=(K_1,\dots,K_\ell)$, $\mathcal Q(\mathcal K)$, $O$, $\operatorname{key}_{\mathcal K}$ & the list of upsets of a node, its qualifier set, the old high set $W\cap\mathcal Q(\mathcal K)$, the key & Section 4, (4.1) \\
(I1), (I2), (I3) & the three invariants of a node & Section 4.1 \\
$X$, $Y$ & $W\setminus D_a$, $W\setminus P_b$ & Lemma 4.2 \\
center $v$; same, switched side; working frame; $E$ & the marked separator of a node; its two sides; the frame of a side; the open part between $a$ and $v$ & Section 4.2 \\
$S_*$; $S_0\subseteq\dots\subseteq S_q$ & the cutoff separated by the center; the cutoffs of a side & (4.3), Lemma 4.3 \\
$s_0\leq\dots\leq s_q<v$ & the positions of the side walk & (4.4) \\
$W'$; $a'$, $b'$, $S$ & a child; its edge data & (4.5) \\
carry, injection $I$, pruning, switch pair, $G$, $h$ & the operations that make the list of a child & Section 5, (5.1) to (5.4), (6.8) \\
run; initial entries & a chain of same-side updates; the entries it starts with & Section 5.2 \\
$\Phi$, $\Phi'$; $H$, $H^*$ & a carry union and its substitute on $U$; the low cutoff of an earlier edge and its extension to $J$ & Lemma 6.2 \\
$a_-$, $b_-$, $v_-$ & boundaries and center of the parent, in the switched frame & Lemma 6.3 \\
$A_m$, $K_m$ & $6m+9$, $12m+25$ & (6.9), (8.11) \\
$\widetilde S_i$ & the cutoff predicate, $\widetilde S_i\cap W=S_i$ & (8.1), (8.2) \\
$B$, $\mathcal{S}$ (Section 9 only) & the bound on the number of formulas; the number of candidate states & (2.1), (9.2) \\
\end{longtable}}

\noindent The numeric state bound is written $\mathcal{S}$ in Section 9; it is the quantity denoted $S$ there in the supplied return. Cutoffs retain $S$. The set variables $L,Z,R$ in state triples are local to those definitions; $L$ in complexity bounds is the encoded input length. The leaf budget $K$ and the indexed list entries $K_j$ are distinguished by their definitions.


\setcounter{section}{1}

\section{Foundations, descriptions, and the finite search}

\textbf{Conventions.} All job sets are finite. Maps between job sets, integer intervals, keys, bit vectors and formulas below are ordinary functions of sets; order-preservation is asserted and proved when needed. All restrictions of precedence use the transitive order on the full job set, not merely the arcs induced by a smaller set.

\begin{numbered}{Lemma}{2.1}{reachability and a rank order}

Let \(G=(V,E)\) be a finite directed acyclic graph. Define \(x\prec y\) if there is a directed path of positive length from \(x\) to \(y\).

\begin{enumerate}
\def\labelenumi{\arabic{enumi}.}
\tightlist
\item
  The relation \(\prec\) is irreflexive and transitive.
\item
  For a function \(\tau:V\to\mathbb Z\), respecting every arc strictly is equivalent to respecting every pair \(x\prec y\) strictly.
\item
  Repeatedly removing a vertex with no remaining predecessor produces a total order extending \(\prec\). With the smallest identifier as a tie-breaker this order is deterministic.
\end{enumerate}

\end{numbered}
\uses{definitions only.}

\begin{proof}
\step{1}{reachability is a strict order}
Transitivity follows by concatenating paths; a positive-length path from a vertex to itself would be a directed cycle.
\by{acyclicity and the path definition}

\step{2}{arc and reachability inequalities agree}
Strict inequalities along a path imply the inequality between its endpoints, and every arc is a path.
\by{the path definition; transitivity of strict integer inequalities}

\step{3}{a deterministic extending order exists}
For the last claim, if a finite nonempty remaining set had no minimal element, following predecessors within it would eventually repeat a vertex and give a cycle. Thus a minimal element always exists. If \(x\prec y\), vertex \(y\) cannot be removed while \(x\) remains. The removal order therefore extends \(\prec\).
\by{finite acyclicity; the least-identifier rule}
\end{proof}

The reachability relation can be computed by a search from each vertex. For a padded set \(J\) of \(N\) jobs, write \[
\lambda:J\longrightarrow[N]=\{1,\ldots,N\}
\] for the bijection sending a job to its position in this deterministic total order. It is strictly order-preserving: \(x\prec y\) implies \(\lambda(x)<\lambda(y)\).

\begin{numbered}{Lemma}{2.2}{integer start times}

Suppose a finite collection of unit jobs has nonnegative real start times and a machine-assignment function \(\mu:V\to[m]\), with no overlap on a machine and with every precedence completed before its successor starts. There is a feasible integer-slot schedule whose makespan is no larger.

\end{numbered}
\uses{Lemma 2.1; definitions of a unit-job schedule and machine assignment.}

\begin{proof}
\step{1}{the augmented graph is acyclic}
On each machine, order the jobs by their original start times and add an arc from each job to the next. Each original or added arc increases the original start time by at least one, so the augmented graph is acyclic.
\by{unit processing times; feasibility of the original schedule}

\step{2}{the recurrence gives integer starts no later than the old starts}
In a topological order of it, define \[
s'(v)=\max\bigl(\{0\}\cup\{s'(u)+1:u\text{ is an augmented immediate predecessor of }v\}\bigr).
\] This defines a function \(s':V\to\mathbb Z_{\geq0}\). For the induction on position in the topological order, a vertex with no incoming arc has \(s'=0\leq s\); assume \(s'(u)\leq s(u)\) for earlier vertices. The recurrence then gives \(s'(v)\leq s(v)\), since every term in the maximum is at most the original \(s(v)\).
\by{Lemma 2.1; induction in the augmented topological order}

\step{3}{the new schedule is feasible}
Added machine-order arcs prevent overlap; original arcs enforce precedence. Setting \(\tau(v)=s'(v)+1\) gives the required slot schedule, with completion time at most the original completion time for each job. Nothing depends on \(m=3\).
\by{the augmented arcs and the recurrence}
\end{proof}

\begin{numbered}{Lemma}{2.3}{padding}

For a deadline \(1\leq T\leq n\), reject if \(n>mT\). Otherwise add \(mT-n\) distinct isolated jobs. Let \(J\) be the resulting set and \(N=mT\leq mn\). A \textbf{full schedule} of \(W\subseteq J\) is a finite sequence of \(m\)-element antichains partitioning \(W\), with \(x\prec y\) implying that the slot of \(x\) occurs strictly earlier than that of \(y\). The empty sequence is a full schedule of the empty set. The original deadline is feasible if and only if \(J\) has a full schedule.

\end{numbered}
\uses{definitions only.}

\begin{proof}
\step{1}{padding gives a full schedule}
A deadline-feasible original schedule leaves exactly \(mT-n\) positions among the \(m\) machines and \(T\) slots. Fill those positions with the distinct isolated jobs.
\by{the deadline schedule and isolated-job definition}

\step{2}{restriction gives the original deadline schedule}
Conversely, deleting isolated padding jobs from a full schedule leaves at most \(m\) original jobs in each slot and preserves all original inequalities. Since \(N\leq mn\), for fixed \(m\) the padded instance has polynomial size.
\by{full-schedule definition; $T\leq n$}
\end{proof}

\begin{numbered}{Lemma}{2.4}{open-interval replacement}

In a full schedule of \(J\), let \(W\) consist of all jobs in the slots strictly between two boundary positions \(a<b\), where a boundary can be the virtual position before the schedule or after it. Replacing these slots by any full schedule of \(W\) preserves a full schedule of \(J\).

\end{numbered}
\uses{definitions only.}

\begin{proof}
\step{1}{the replacement occupies the same interval}
Every full schedule of \(W\) has exactly \(|W|/m\) slots, so it fits the same positions.
\by{the definition of a full schedule}

\step{2}{every precedence pair is respected}
Relations with both endpoints in \(W\) are respected by the replacement, and those with neither endpoint in \(W\) are unchanged. An exterior job has a slot before the entire interval or after the entire interval. Its temporal relation with every interior job therefore remains the same. This verifies every relation with precisely one endpoint in \(W\).
\by{the original schedule; the open-interval definition}
\end{proof}

The replacement ranges over full schedules of the interval job set, each respecting its restricted transitive precedence order.

\subsection{Frames, cones, and atomic formulas}

There are two \textbf{frames}. The forward frame uses the given precedence, time and rank. The reverse frame uses \[
x\prec_- y\iff y\prec x,\qquad
\lambda_-(x)=N+1-\lambda(x).
\] For a schedule of \(T\) slots its time-reversal map is the bijection \([T]\to[T]\), \(t\mapsto T+1-t\). This reverses the sequence of slots and changes a full schedule for \(\prec\) into one for \(\prec_-\). Applying it twice is the identity. All terms such as ``before'', ``upset'', ``rank'' and ``predecessor'' in an argument refer to its stated frame.

For any \(m\)-element subset \(a\subseteq J\), define \[
P_a=\{x\in J:(\exists y\in a)\ x\prec y\},\qquad
D_a=\{x\in J:(\exists y\in a)\ y\prec x\},
\] and \(\widehat P_a=P_a\cup a\), \(\widehat D_a=D_a\cup a\). These are defined for every \(m\)-element subset, although the construction uses them at antichain slots. A subset \(K\subseteq J\) is an \textbf{upset} if \(x\in K\) and \(x\prec y\) imply \(y\in K\); a \textbf{downset} satisfies the reversed implication. By transitivity, \(D_a,\widehat D_a\) are upsets and \(P_a,\widehat P_a\) are downsets. Complements interchange upsets and downsets; unions and intersections preserve either class. Rank suffixes are upsets and rank prefixes are downsets. Reversal exchanges \(P\) with \(D\), including weak cones.

For the virtual boundary before the schedule set \[
D=\widehat D=J,\qquad P=\widehat P=\varnothing;
\] for the virtual boundary after it set \[
D=\widehat D=\varnothing,\qquad P=\widehat P=J.
\] The virtual boundaries contain no jobs.

\begin{numbered}{Lemma}{2.5}{cones and time}

At a genuine slot \(a\) of a full schedule, jobs of \(P_a\) occur strictly before \(a\), jobs of \(D_a\) strictly after \(a\), and no job of \(a\) belongs to either strict cone. In particular, no job at or after \(a\) belongs to \(P_a\), and no job at or before \(a\) belongs to \(D_a\).

\end{numbered}
\uses{definitions only.}

\begin{proof}
A witness \(x\prec y\) with \(y\in a\) forces the slot of \(x\) to precede \(a\). The successor statement is identical with directions reversed. Members of \(a\) are pairwise incomparable, and no job precedes itself. The virtual-boundary conventions satisfy the analogous temporal implications.
\end{proof}

A \textbf{formula} is a finite full binary tree whose internal vertices are labelled intersection or union, and whose leaves are one of the following subsets or the complement of one of them:

\begin{itemize}
\tightlist
\item
  \(J\) and \(\varnothing\);
\item
  a rank threshold \(\{\lambda\leq k\}\) or \(\{\lambda\geq k\}\), with \(0\leq k\leq N+1\), in either frame;
\item
  \(a,P_a,D_a,\widehat P_a,\widehat D_a\) for an \(m\)-element subset \(a\), in either frame.
\end{itemize}

Its size is its number of leaves. Evaluation is the function from the set of such formula trees to \(2^J\) obtained by interpreting the leaves and internal operations. Write \([F]\) for the evaluated set. For \(K\geq1\) let \[
\mathcal F_K=\{[F]:F\text{ has at most }K\text{ leaves}\}.
\] A global description of \(W\) means \(W\in\mathcal F_K\). A relative description of \(L\subseteq W\) means \(L=W\cap[F]\) for some such \(F\); it does not require the intersection with \(W\) to have a short global description.

\begin{numbered}{Lemma}{2.6}{complements}

Every formula has a complement formula of exactly the same size.

\end{numbered}
\uses{definitions only.}

\begin{proof}
\step{1}{define the transformation and verify its leaf base case}
Complement each leaf and interchange union with intersection at every internal vertex. Induct on the number of vertices of the formula tree. At a leaf, the allowed complement flag gives exactly the complementary set, and both formulas have one leaf.
\by{the permitted complement flags}

\step{2}{the induction step preserves complements and leaf counts}
At an internal vertex, assume that the transformed child formulas give the complements of their old evaluations and preserve their leaf counts. The corresponding De Morgan identity then gives the complementary evaluation at this vertex. Both trees have the sum of their child leaf counts. Thus the new evaluation is \(J\setminus[F]\), with no leaf created or deleted.
\by{the induction hypothesis; the two De Morgan identities}
\end{proof}

Thus set differences used below are formulas, with leaf counts equal to the sum of the counts of the two operands.

\begin{numbered}{Lemma}{2.7}{finite formula enumeration}

For fixed \(m\geq3\) and \(K\), the number of formulas with at most \(K\) leaves is at most \[
B=C_{m,K}N^{mK}.
\tag{2.1}
\] Their evaluated sets can be enumerated and deduplicated as bit vectors by a deterministic algorithm polynomial in \(N\).

\end{numbered}
\uses{Lemma 2.1; definitions of the atomic alphabet and formula evaluation.}

\begin{proof}
\step{1}{bound the atomic alphabet}
There are \(\binom Nm\leq N^m\) slot subsets. Counting both frames and complements, the leaf alphabet has at most \[4+8(N+2)+20\binom Nm\leq64N^m\] elements for \(N\geq m\geq3\): the constants contribute four choices, the thresholds contribute \(8(N+2)\), and the five subset/cone types contribute \(20\binom Nm\). The inequality follows from \(8N+20\leq16N\leq16N^m\) and \(\binom Nm\leq N^m\).
\by{the permitted atoms; $N\geq m\geq3$}

\step{2}{count the formula trees}
For \(\ell\) leaves there are at most \(4^\ell\) full binary tree shapes, at most \(2^{\ell-1}\) internal labellings, and at most \((64N^m)^\ell\) leaf labellings. Summing gives, for example, \(C_{m,K}=K\,512^K\).
\by{the binary-tree shape and label counts}

\step{3}{enumerate and retain every evaluated set}
Enumerate these finite choices, evaluate each using the reachability table, and compare evaluated \(N\)-bit vectors in pairs to remove duplicates. For fixed \(m,K\), every loop has polynomially many iterations.
\by{Lemma 2.1; finite formula evaluation; bit-vector equality}
\end{proof}

\subsection{The dynamic program}

Process the sets of \(\mathcal F_K\) in increasing cardinality. A boolean function \[
\mathsf A:2^J\longrightarrow\{\mathrm{false},\mathrm{true}\}
\] is initialized as false outside \(\mathcal F_K\), and \(\mathsf A(\varnothing)=\mathrm{true}\). A set whose size is not divisible by \(m\) is false.

For nonempty \(W\in\mathcal F_K\) with \(m\mid |W|\), a \textbf{state} is a disjoint decomposition \[
W=L\mathbin{\dot\cup}Z\mathbin{\dot\cup}R
\] such that \(Z\) is an \(m\)-element antichain, \(L\) has a relative description in \(W\), and the following compatibility tests hold:

\begin{itemize}
\tightlist
\item
  no job in \(Z\cup R\) precedes a job in \(L\);
\item
  no job in \(R\) precedes a job in \(Z\).
\end{itemize}

The state is initial when \(\mathsf A(L)\) is true and final when \(\mathsf A(R)\) is true. There is a directed transition to \((L',Z',R')\) exactly when \[
L\cup Z\subseteq L',
\qquad
\mathsf A\bigl(L'\setminus(L\cup Z)\bigr)=\mathrm{true}.
\tag{2.2}
\] Declare \(W\) accepted when an initial state reaches a final state, allowing a path with one state. Every acceptance lookup in this computation concerns a smaller set: \(L\) and \(R\) omit \(Z\), while a transition gap omits at least the nonempty \(Z\). Thus the procedure is well-defined.

\begin{numbered}{Proposition}{2.8}{soundness and witnesses}

For any \(m\geq1\) and any description family used in this procedure, an accepted set \(W\) has a full schedule. Such a schedule can be reconstructed from the accepted states.

\end{numbered}
\uses{definitions of the state graph, compatibility, and acceptance in Section 2.2; induction on set cardinality.}

\begin{proof}
\step{1}{form the ordered partition from an accepting path}
Induct on \(|W|\). The base \(W=\varnothing\) has the empty schedule. Assume that every accepted smaller set has a full schedule. The empty case is immediate. For a path of states \((L_i,Z_i,R_i)\), \(1\leq i\leq k\), from initial to final, put \[
A_0=L_1,\quad
A_i=L_{i+1}\setminus(L_i\cup Z_i)\ (1\leq i<k),\quad
A_k=R_k.
\] The inclusion in (2.2), together with disjointness of each state, gives the ordered disjoint partition \[
W=A_0\mathbin{\dot\cup}Z_1\mathbin{\dot\cup}A_1
\mathbin{\dot\cup}\cdots\mathbin{\dot\cup}Z_k\mathbin{\dot\cup}A_k.
\tag{2.3}
\] For completeness of this set identity, induction on \(i\) gives \(L_{i+1}=A_0\dot\cup Z_1\dot\cup A_1\dot\cup\cdots\dot\cup Z_i\dot\cup A_i\); adding the last state's \(Z_k,R_k\) proves (2.3).
\by{the induction hypothesis; transition inclusion (2.2)}

\step{2}{concatenation respects every precedence}
Each \(A_i\) is accepted and smaller than \(W\), so it has a full schedule by induction. Concatenate these schedules and the individual slots \(Z_i\) in the displayed order. A backward precedence from \(Z_j\) into an earlier piece would run from \(Z_j\) into \(L_j\), violating its state's compatibility. A backward precedence from \(A_j\), \(j\geq1\), into an earlier piece would run from \(R_j\) into \(L_j\cup Z_j\), again violating compatibility. These cover every relation between different pieces. Relations inside pieces are respected by their schedules. The concatenation is full.
\by{the induction hypothesis; state compatibility}

\step{3}{decode a witness from the stored path}
Store the predecessor of each reached state and the schedule of each accepted smaller set. Backtracking from a reached final state yields the path and hence the displayed schedule. This uses only one \(m\)-element slot in each state.
\by{the partition (2.3) and the accepted smaller witnesses}
\end{proof}

\begin{numberedup}{Definition}{2.9}{description hierarchy}

A hierarchy consists of nonempty job sets organized in a rooted tree, together with full-schedule witnesses, as follows.

\begin{itemize}
\tightlist
\item
  The root is \(J\) between the two virtual boundaries.
\item
  A node \(W\) is the job set of an open interval in its witness schedule of \(J\). Its interval may be rescheduled, and then at least one slot is marked.
\item
  Its children are exactly the nonempty gaps between consecutive marks or between a boundary and its nearest mark. Each child receives a copy of the parent's current witness schedule.
\item
  Every node has a global description using at most \(K\) leaves. For each mark \(z\), the set of jobs of \(W\) before \(z\) in original time has a relative description using at most \(K\) leaves.
\end{itemize}

Copies separate the witnesses for distinct branches; they are not algorithmic state variables. Every child omits at least one marked slot, so the tree has finite depth.

\end{numberedup}
\uses{definitions of full schedules, open intervals, and global and relative descriptions.}

\begin{numbered}{Lemma}{2.10}{hierarchy implies acceptance}

If \(J\) has such a hierarchy, the dynamic program accepts \(J\).

\end{numbered}
\uses{Definition 2.9; state and acceptance definitions in Section 2.2.}

\begin{proof}
\step{1}{each mark gives a candidate state}
Induct on the number of slots \(|W|/m\) of a node. For one slot, that slot is marked and all gaps are empty, so it is an initial and final state. For a larger node assume the claim for its smaller children. List its marked slots in original time as \(z_1,\ldots,z_k\). Let \(L_i\) and \(R_i\) be its jobs before and after \(z_i\). Each \((L_i,z_i,R_i)\) is a state: the relative description is part of the hierarchy, and compatibility follows from the witness schedule.
\by{Definition 2.9; the state definition}

\step{2}{the state path accepts the node}
The first \(L_1\), the last \(R_k\), and each gap \(L_{i+1}\setminus(L_i\cup z_i)\) are empty or are smaller children. They are accepted by induction. Hence the first state is initial, the last final, and successive states are connected by (2.2). The node is accepted. Apply this at the root.
\by{the size-induction hypothesis; transition definition (2.2)}
\end{proof}

The rest of the structural proof constructs a hierarchy with \(K=K_m=12m+25\).

\section{Separators and the shared-cutoff walk}

Fix one frame. A \textbf{domain} \(U\) is the job set of an open interval of a full schedule of \(J\). A \textbf{cutoff} \(S\subseteq U\) is a downset relative to \(U\): if \(x\in S\), \(y\in U\) and \(y\prec x\), then \(y\in S\). Jobs in \(S\) are called low and those in \(U\setminus S\) high.

A slot \(z\) in the domain is a \textbf{separator at \(S\)} when

\begin{enumerate}
\def\labelenumi{(\alph{enumi})}
\item
  every high job before \(z\) belongs to \(P_z\); and
\item
  every low job after \(z\) belongs to \(D_z\).
\end{enumerate}

Nothing is required of the jobs at \(z\) itself. The boundary \(a\) immediately before the domain is a separator if all low jobs of \(U\) belong to \(D_a\); its condition (a) is empty. Put \[
B_z(H)=P_z\cup(H\setminus\widehat D_z)
\quad(H\subseteq J).
\tag{3.1}
\]

\begin{numbered}{Lemma}{3.1}{separator descriptions}

If \(z\) is a slot of \(U\) separating \(S\), and \(H\subseteq J\) satisfies \(H\cap U=S\), then \(B_z(H)\cap U\) is exactly the set of jobs of \(U\) before \(z\).

\end{numbered}
\uses{Lemma 2.5; separator definition and (3.1).}

\begin{proof}
\step{1}{every prefix job satisfies the formula}
A high job before \(z\) belongs to \(P_z\) by (a). A low job before \(z\) belongs to \(H\) and, by Lemma 2.5, not to \(\widehat D_z\). This proves inclusion of the prefix.
\by{separator clauses; Lemma 2.5}

\step{2}{every selected job belongs to the prefix}
Conversely, a job of \(P_z\) is before \(z\). A job of \(U\cap H\setminus\widehat D_z\) is low, is not in slot \(z\), and cannot be after \(z\), because condition (b) would put it in \(D_z\). Thus it too is before \(z\).
\by{Lemma 2.5; separator clause (b)}
\end{proof}

\begin{numbered}{Lemma}{3.2}{reversal of a separator}

A slot \(z\) of \(U\) separates \(S\) in one frame if and only if it separates \(U\setminus S\) in the reverse frame.

\end{numbered}
\uses{definitions of frame reversal, cutoff, and separator.}

\begin{proof}
\step{1}{reversal sends the cutoff to its complement}
The complement of a relative downset is a relative upset, hence a cutoff in the reversed order.
\by{the definitions of a cutoff and the reversed order}

\step{2}{the separator clauses are exchanged}
In the reverse frame, high jobs are the original low jobs, ``before'' means original ``after'', and \(P_z\) is original \(D_z\). Thus the reversed condition (a) is exactly original (b). Similarly reversed (b) is original (a). Reversing again proves both directions.
\by{the time and cone reversal maps}
\end{proof}

\subsection{Normalization}

A \textbf{priority order} on \(U\) is a total order extending its precedence restriction. Let \(\pi:U\to[|U|]\) be the bijection giving its positions, smaller meaning better priority. Encode a full schedule of \(U\) by sorting each slot increasingly by \(\pi\) and concatenating the resulting \(m\)-tuples in time order. A \textbf{normalized} full schedule minimizes this finite integer sequence lexicographically among all full schedules of \(U\). Such a schedule exists whenever \(U\) is schedulable. Normalized schedules are existence witnesses, not schedules that the finite search must compute.

\begin{numbered}{Lemma}{3.3}{no improving exchange, including suffix rearrangements}

Let a normalized schedule of \(U\) and one of its slots \(r\) be fixed. Replace the slots strictly after \(r\) by any full schedule of the same suffix job set, keeping a full schedule of \(U\). A job \(y\) at \(r\) cannot then be exchanged with a later job \(x\) satisfying \(\pi(x)<\pi(y)\) so that a full schedule of \(U\) results.

\end{numbered}
\uses{the definition of a normalized schedule (Section 3.1).}

\begin{proof}
Suppose such an exchange of $y$ with $x$ gives a full schedule $\Sigma'$ of $U$. Write $\Sigma$ for the normalized
schedule.
\step{1}{$\Sigma'$ and $\Sigma$ agree before $r$} The replacement changes only slots strictly after $r$, and the
exchange changes only the slot $r$ and one later slot. \by{hypotheses of the lemma}
\step{2}{the sorted $m$-tuple of slot $r$ decreases lexicographically} At $r$ the set of priority values changes by
replacing $\pi(y)$ with the smaller value $\pi(x)$, which is distinct from all values at $r$. For every threshold $q$
the number of values at most $q$ does not decrease, so every order statistic of the $m$ values weakly decreases. The
sum of the $m$ values decreases strictly, so some order statistic decreases strictly. Hence the first coordinate in
which the two sorted $m$-tuples differ is smaller in $\Sigma'$. \by{$\pi$ is a bijection; $\pi(x)<\pi(y)$}
\step{3}{contradiction} By Steps 1 and 2 the encoded sequence of $\Sigma'$ is lexicographically smaller than that of
$\Sigma$, against the minimality that defines a normalized schedule. \by{definition of normalized, Section 3.1}

\noindent\emph{Where $m$ enters:} nowhere beyond the size of a slot. The argument holds for every $m\geq 1$ and uses
no exchange involving the other $m-1$ jobs of the slot.
\end{proof}

An exchange in all following proofs swaps two jobs between their slots. It preserves the size of both slots. Feasibility will always be checked through predecessors and successors of those two jobs; all other pairs retain their order.

\begin{numbered}{Lemma}{3.4}{one advance}

Let \(U\) have a normalized full schedule for a priority order, and let \(M\subseteq L\) be two prefixes of this order. A position is either a slot of \(U\) or the boundary \(a\) before \(U\). Suppose position \(s\) satisfies separator condition (a) for \(M\).

An \textbf{exception for \(L\) at \(s\)} is a job \(x\in L\) after \(s\) with \(x\notin D_s\). Suppose exceptions exist, let \(t\) be their earliest slot, and choose an exception \(x\) at \(t\). If \(x\in M\), then:

\begin{enumerate}
\def\labelenumi{\arabic{enumi}.}
\tightlist
\item
  every predecessor of \(x\) belonging to \(U\) is strictly before \(s\);
\item
  for every \(y\in U\setminus M\) at a slot \(r\) of \(U\) with \(s\leq r<t\), some job of \(U\setminus M\) at \(t\) is a strict successor of \(y\);
\item
  \(t\) satisfies condition (a) for \(M\).
\end{enumerate}

\end{numbered}
\uses{Lemma 3.3; definitions of domain, priority prefix, exception, and separator.}

\begin{proof}
\step{1}{every predecessor of the exception is before the old position}
A predecessor \(w\in U\) of \(x\) belongs to \(L\) because \(L\) is a prefix of an order extending precedence. It is before \(t\). If it were in \(s\), then \(x\in D_s\). If it were strictly between \(s\) and \(t\), the minimality of \(t\) would give \(w\in D_s\) and hence \(x\in D_s\) by transitivity. Both are contradictions. This proves (1), including the case when \(s\) is exterior and contains no job of \(U\).
\by{priority-prefix closure; earliest-exception choice; transitivity}

\step{2}{a forbidden exchange forces a successor}
For (2), suppose \(y\) has no successor in \(U\) at or before \(t\). Exchange \(x\) and \(y\). All predecessors of \(x\) in \(U\) are before \(s\leq r\) by (1), and all its successors were after \(t\). Thus moving \(x\) to \(r\) violates no relation involving \(x\). The predecessors of \(y\) were before \(r\) and its successors in \(U\) are after \(t\) by the supposition, so moving \(y\) to \(t\) violates none involving \(y\). In particular, there is no forbidden relation with another job occupying either new slot. Exterior predecessors of an interior job are before the whole interval, and exterior successors are after it, by the original full schedule. Thus the exchange is feasible in \(U\) and also in its full-schedule witness.

But \(x\in M\) and \(y\notin M\), so \(x\) has better priority. Lemma 3.3 forbids the exchange. Therefore \(y\) has a successor \(y_1\in U\) at some time \(r<r_1\leq t\).
\by{Step 1; Lemma 3.3}

\step{3}{successor propagation reaches the target slot}
This successor is outside \(M\), since \(M\) is a cutoff. If \(r_1<t\), apply the same argument to \(y_1\). The times strictly increase in a finite interval. The process reaches a successor at \(t\), still outside \(M\). This proves (2).
\by{prefix closure; strictly increasing successor times; Step 2}

\step{4}{the target satisfies the high-job separator clause}
For (3), let \(y\in U\setminus M\) be before \(t\). If its slot is at or after \(s\), use (2). Otherwise \(s\) must be a slot of \(U\), and condition (a) at \(s\) gives a job \(w\) in \(s\) with \(y\prec w\). This \(w\) is outside \(M\), so (2) gives \(w\prec w'\) for a job \(w'\) at \(t\). Hence \(y\prec w'\) and \(y\in P_t\).
\by{Step 3; the old separator clause (a); transitivity}
\end{proof}

\begin{numbered}{Lemma}{3.5}{fixed-cutoff walk}

Let \(L\) be a priority prefix and suppose \(s\) satisfies (a) for \(L\). Repeatedly move from the current position to the earliest slot with an exception for \(L\), stopping when none exists. The process terminates at a separator for \(L\), never moves backward, and does not move from a position already separating \(L\).

\end{numbered}
\uses{Lemma 3.4.}

\begin{proof}
Lemma 3.4 with \(M=L\) preserves (a) at each target. Every move strictly increases the position. There are only finitely many positions. At termination the absence of exceptions is exactly (b). If the starting position already separates \(L\), there is no first exception.
\end{proof}

\begin{numbered}{Lemma}{3.6}{shared-cutoff walk}

Let \(S\) and \(S^+=S\cup\{p\}\) be consecutive priority prefixes. From a separator \(s\) for \(S\), run the fixed-cutoff walk for \(L=S^+\). Its final position \(t\geq s\) separates both \(S\) and \(S^+\).

\end{numbered}
\uses{Lemmas 3.4 and 3.5.}

\begin{proof}
\step{1}{the no-advance case already separates both cutoffs}
Condition (a) for \(S^+\) already holds at \(s\), because the high set shrinks. If there is no exception for \(S^+\), then \(s\) separates both and there is no move.
\by{the old separator and $S\subseteq S^+$}

\step{2}{the first target separates the high jobs of both cutoffs}
Otherwise every old low job after \(s\) belongs to \(D_s\), so the only possible first exception is \(p\). The first target \(t_1\) is the slot of \(p\). By Lemma 3.4 with \(M=L=S^+\), it satisfies (a) for \(S^+\). It also satisfies (a) for \(S\): the only difference between the two high sets is \(p\), which is at \(t_1\), not before it.
\by{Lemma 3.4; the single promoted job}

\step{3}{later advances preserve both clauses and terminate}
At each subsequent move, the exception is after \(t_1\), and hence is a job of \(S^+\) other than \(p\), that is, a job of \(S\). Apply Lemma 3.4 once with \(M=S\) and once with \(M=S^+\) to preserve (a) for both cutoffs. At the final position, no exception for \(S^+\) remains. This proves (b) for \(S^+\) and for its subset \(S\).
\by{Lemmas 3.4 and 3.5; $S\subseteq S^+$}
\end{proof}

The same result applies if the priority order and prefixes are on a larger set \(W\supseteq U\), with normalization on the restriction to \(U\). The intersection of a priority prefix with \(U\) is a prefix of the restricted order. Promoting a job outside \(U\) leaves that intersection unchanged, so the walk does not move.

\begin{numbered}{Corollary}{3.7}{a separator can be created}

Let nonempty \(U\) be a domain with a full schedule, and let \(S\) be any cutoff in it. Some full schedule of \(U\) has a slot separating \(S\).

\end{numbered}
\uses{Lemma 3.6; definitions of cutoff and normalization.}

\begin{proof}
\step{1}{the chosen priority order respects precedence}
Order \(S\) first, then \(U\setminus S\), each by rank. A relation from the second group into the first would contradict that \(S\) is a cutoff, so this is a priority order.
\by{the cutoff property and rank order}

\step{2}{the promotions produce a separator slot}
Normalize a full schedule for it. The first slot separates the empty cutoff. Promote the members of \(S\) in that order, applying Lemma 3.6 each time. The final position is a slot and separates \(S\).
\by{normalization; Lemma 3.6}
\end{proof}

\section{Centers, sides, and the hierarchy data}

Every node \(W\) has a full-schedule witness, a tag consisting of one frame, and boundaries \(a,b\) in that frame. It carries an ordered list \[
\mathcal K=(K_1,\ldots,K_\ell)
\] of upsets of the whole set \(J\) in its tag frame. Its \textbf{qualifier set} is \[
\mathcal Q(\mathcal K)=\bigcup_{j=1}^{\ell}K_j,
\] and its \textbf{old high set} is \(O=W\cap\mathcal Q(\mathcal K)\). On \(\mathcal Q(\mathcal K)\) define the function \[
\operatorname{key}_{\mathcal K}:
\mathcal Q(\mathcal K)\longrightarrow[\ell]\times[N],
\qquad
x\longmapsto\bigl(\max\{j:x\in K_j\},\lambda(x)\bigr),
\tag{4.1}
\] with lexicographic order in the target. For an empty list there are no qualifiers and no keys. Keys of distinct qualifiers are distinct, because \(\lambda\) is injective.

Each \(K_j\) is a fixed subset given by a formula in fixed job subsets and rank thresholds. Changing a witness schedule later does not change that subset.

\begin{numbered}{Lemma}{4.1}{keys respect precedence}

If \(x\) qualifies and \(x\prec y\), then \(y\) qualifies and \(\operatorname{key}_{\mathcal K}(x)<\operatorname{key}_{\mathcal K}(y)\).

\end{numbered}
\uses{definitions only.}

\begin{proof}
Every upset entry containing \(x\) contains \(y\). The largest entry index of \(y\) is therefore at least that of \(x\). If it is larger, the first coordinate proves the claim; if equal, \(\lambda(x)<\lambda(y)\) proves it.
\end{proof}

\subsection{The three invariants}

Every node is required to satisfy the following statements in its tag frame.

\textbf{(I1)} \(O\subseteq P_b\) and \(W\setminus O\subseteq D_a\).

\textbf{(I2)} Every qualifier scheduled at or before \(a\) belongs to \(\widehat P_a\).

\textbf{(I3)} After replacing the interval of \(W\) by \emph{any} full schedule of \(W\), no exchange of a job \(x\in W\setminus D_a\) with a qualifier \(y\) in slot \(a\), satisfying \(\operatorname{key}_{\mathcal K}(x)<\operatorname{key}_{\mathcal K}(y)\), gives a full schedule of \(J\).

The key of \(x\) in (I3) is defined: (I1) implies \(W\setminus D_a\subseteq O\). When \(a\) is virtual, (I2) and (I3) are empty assertions.

At the root, use the forward frame, virtual boundaries, and \(\mathcal K=(J)\). Then \(O=J\), (I1) holds by the virtual conventions, and the other two invariants are vacuous.

\begin{numbered}{Lemma}{4.2}{cover and two cutoffs}

In either frame, writing \(a,b\) for the node's boundaries in that frame and \[
X=W\setminus D_a,\qquad Y=W\setminus P_b,
\] one has \[
W\subseteq D_a\cup P_b,\qquad X\subseteq P_b,\qquad Y\subseteq D_a.
\tag{4.2}
\] Both \(X\) and \(W\cap P_b\) are cutoffs in \(W\).

\end{numbered}
\uses{invariant (I1); frame and cone definitions.}

\begin{proof}
\step{1}{the cover and inclusions hold in the tag frame}
In the tag frame, (I1) puts the two parts \(O,W\setminus O\) in \(P_b,D_a\) respectively. Thus a job outside \(D_a\) is in \(O\subseteq P_b\), and a job outside \(P_b\) is in \(W\setminus O\subseteq D_a\).
\by{invariant (I1)}

\step{2}{the same assertions hold in the reversed frame}
Reversal interchanges the boundary names and the predecessor/successor cones; it exchanges the roles of \(X,Y\) and leaves the assertions equivalent.
\by{the explicit time, boundary, and cone reversal maps}

\step{3}{the two sets are relative downsets}
Finally \(J\setminus D_a\) and \(P_b\) are downsets, and their intersections with \(W\) are relative downsets.
\by{global cone closure and taking intersections with $W$}
\end{proof}

\subsection{Choosing the center and treating the two sides}

In the tag frame, use Corollary 3.7 to reschedule \(W\) so that a slot \(v\) separates \(W\cap P_b\) on the whole domain \(W\). Mark this \textbf{center}. At the root use its last slot as the center: with cutoff \(J\), (a) has no high job and (b) has no later job.

Treat each side of \(v\) as the part before \(v\) in a \textbf{working frame}. The side before \(v\) in the tag frame is the \textbf{same side}; the side after it uses the reversed frame and is the \textbf{switched side}. On either side let \(a,b\) be the node boundaries in that working frame and let \(E\) be the open part between \(a\) and \(v\).

The center separates, on \(W\), the cutoff \[
S_*=
\begin{cases}
W\cap P_b,&\text{same side},\\
X=W\setminus D_a,&\text{switched side}.
\end{cases}
\tag{4.3}
\] The second assertion follows from Lemma 3.2: the new cutoff is the complement of the old one in \(W\), with old \(P_b\) becoming new \(D_a\).

On a same side set \(S_0=(W\setminus O)\cap P_b\), order \(S_0\) by rank, then \(O\) by key, then \(Y\) by rank; promote all jobs of \(O\) in key order. On a switched side start with \(S_0=\varnothing\), order \(X\) by rank and then \(W\cap D_a\) by rank; promote all jobs of \(X\) in rank order. These orders are on all of \(W\).

\begin{numbered}{Lemma}{4.3}{priority orders and cutoff containment}

Both orders extend precedence. The successive low sets are prefixes \[
S_0\subseteq S_1\subseteq\cdots\subseteq S_q=S_*,
\] each addition being one promoted job. Every \(S_i\) is contained in \(P_b\), and \(S_0\subseteq D_a\).

\end{numbered}
\uses{Lemmas 4.1 and 4.2; invariant (I1).}

\begin{proof}
\step{1}{the same-side order and its prefixes have the required properties}
The two working sides, same and switched, exhaust the construction. On the same side, \(S_0,O,Y\) partition \(W\) because \(O\subseteq P_b\). If \(w\prec x\) in \(W\) with \(x\in S_0\), down-closure gives \(w\in P_b\). Also \(w\notin O\): qualification of \(w\) would force qualification of \(x\). Thus \(S_0\) is a cutoff. No job of \(Y=W\setminus P_b\) precedes a job of \(S_0\cup O=W\cap P_b\). Keys extend precedence in \(O\) by Lemma 4.1; ranks do so in the other groups. This proves the priority property and the stated prefixes. Their final union is \(W\cap P_b\). Invariant (I1) gives \(S_0\subseteq D_a\).
\by{Lemmas 4.1 and 4.2; invariant (I1)}

\step{2}{the switched-side order and its prefixes have the required properties}
On the switched side, \(X\) is a cutoff by Lemma 4.2. Ordering it before its complement and using ranks within both groups therefore extends precedence. Its rank prefixes end at \(X\), which lies in \(P_b\) by (4.2), and its initial cutoff is empty.
\by{Lemma 4.2; rank order}
\end{proof}

Normalize the open part \(E\) for the restricted priority order, independently on both sides. Lemma 2.4 preserves the full schedule. The center remains a separator, since the sets of jobs on its two sides do not change.

\subsection{Side walks and children}

The boundary \(a\) separates \(S_0\) on \(E\), since \(S_0\subseteq D_a\). Set \(s_0=a\). For each promotion \(S_{i-1}\to S_i\), apply Lemma 3.6 on \(E\) and let \(s_i\) be the walk's final position. Intermediate positions during a promotion are not necessarily marked. Set \(s_{q+1}=v\). Then \[
a=s_0\leq s_1\leq\cdots\leq s_q<s_{q+1}=v.
\tag{4.4}
\]

\begin{numbered}{Lemma}{4.4}{common cutoffs of adjacent positions}

For \(0\leq i<q\), both \(s_i\) and \(s_{i+1}\) separate \(S_i\) on \(E\). The position \(s_q\) separates \(S_q\) there, and every high job of \(E\) for \(S_q\) belongs to \(P_v\).

\end{numbered}
\uses{Lemmas 3.6 and 4.3; the center and side-walk construction.}

\begin{proof}
\step{1}{successive walk positions share the old cutoff}
The initial position separates \(S_0\). For every promotion, Lemma 3.6 says its final position separates both the previous and the new cutoff. This proves all the paired assertions.
\by{Lemmas 4.3 and 3.6}

\step{2}{the center has the final high-job property}
For the final endpoint \(v\), use that it separates \(S_q=S_*\) on the larger domain \(W\); its condition (a) covers all high jobs of \(E\).
\by{the center-separator property}
\end{proof}

Mark the distinct \(s_i\) that are slots of \(E\), together with \(v\). For each strict step \(s_i<s_{i+1}\), the nonempty gap between them becomes a child \(W'\) with working-frame tag and \textbf{edge data} \[
a'=s_i,\qquad b'=s_{i+1},\qquad S=S_i.
\tag{4.5}
\] Repeated positions do not create a gap. Consecutive distinct positions are separated by exactly one strict step in the nondecreasing sequence (4.4), so these children are exactly the nonempty gaps required by Definition 2.9.

\begin{numbered}{Lemma}{4.5}{edge labels and shrinking}

For each such child, \(b'\) is a slot of \(W\), and \(a'\) is the boundary \(a\) or a slot of \(E\). Moreover \[
W'\cap S\subseteq D_{a'}\cap P_b,\qquad
W'\setminus S\subseteq P_{b'}.
\tag{4.6}
\] The child has fewer jobs than \(W\), and the root has no switched-side child.

\end{numbered}
\uses{Lemmas 4.3 and 4.4; Definition 2.9 and the child construction.}

\begin{proof}
\step{1}{the endpoints and cone inclusions hold}
The position claims follow from (4.4). A low job of the gap is after \(a'\), so it lies in \(D_{a'}\) by the separator property at \(a'\). It lies in \(P_b\) because \(S\subseteq P_b\). A high job is before \(b'\) and belongs to \(P_{b'}\) by Lemma 4.4, including the case \(b'=v\).
\by{the side-walk positions; Lemmas 4.3 and 4.4}

\step{2}{the child shrinks and the root has no switched child}
The gap omits \(v\), giving strict shrinkage. At the root \(v\) is the last slot, so the switched side is empty.
\by{the child and root-center definitions}
\end{proof}

The remaining task is to give each child a short list and to preserve (I1)--(I3). The induction order will be checked explicitly in Section 7.

\section{Carried and pruned lists}

\begin{numbered}{Lemma}{5.1}{suffix carry}

For a nonempty list \(\mathcal K=(K_1,\ldots,K_\ell)\), an index \(1\leq d\leq\ell\), and \(0\leq r\leq N+1\), define \[
\operatorname{carry}_{d,r}(\mathcal K)
=
\bigl(K_d\cap\{\lambda\geq r\},K_{d+1},\ldots,K_\ell\bigr).
\tag{5.1}
\] This is a function, for each permitted \(d,r\), from lists of length at least \(d\) of upsets of \(J\) to lists of upsets of \(J\). Its qualifiers are exactly those old qualifiers whose old key is at least \((d,r)\). On those qualifiers its new key order agrees with the old one.

\end{numbered}
\uses{definitions only.}

\begin{proof}
\step{1}{qualification is exactly the old key suffix}
Let \(j\) be a qualifier's largest old entry index. Exactly one of \(j>d\), \(j=d\), and \(j<d\) holds. If \(j>d\), that entry survives unchanged, so the job qualifies. If \(j=d\), it qualifies precisely when its rank is at least \(r\). If \(j<d\), no kept entry contains it. These three cases give exactly the threshold on the old key.
\by{the three exhaustive cases $j>d$, $j=d$, and $j<d$}

\step{2}{the key-order map preserves comparisons and entries remain upsets}
For every surviving qualifier, its largest qualifying index is decreased by the same amount \(d-1\), and its rank is unchanged. Hence every pairwise comparison is preserved. The first entry remains an upset because it is an intersection of two upsets; the other entries already are upsets.
\by{the index translation $j\mapsto j-d+1$; upset closure}
\end{proof}

\textbf{Pruning on \(W'\)} means deleting every entry disjoint from \(W'\) and retaining the order of the others. It is a function from finite lists of subsets of \(J\) to such lists. It preserves the union's intersection with \(W'\).

\begin{numbered}{Lemma}{5.2}{comparison after pruning}

Suppose a list is pruned on \(W'\). Let \(x\in W'\) and \(y\in J\) both qualify after pruning. If the pruned key of \(x\) is smaller than the pruned key of \(y\), then their old keys have the same strict inequality.

\end{numbered}
\uses{definitions only.}

\begin{proof}
\step{1}{the smaller-key job keeps its largest old qualifying index}
Before renumbering the retained entries, let \(M(z)\) be the largest old index containing \(z\) and \(R(z)\) the largest retained old index containing \(z\). Every entry containing \(x\) meets \(W'\), so \(R(x)=M(x)\); for \(y\) one has \(R(y)\leq M(y)\). Renumbering the retained indices preserves their order. The strict comparison has two cases: \(R(x)<R(y)\), or \(R(x)=R(y)\) and \(\lambda(x)<\lambda(y)\).
\by{pruning; $x\in W'$}

\step{2}{both possible strict key comparisons transfer to the old list}
If \(R(x)<R(y)\), then \(M(x)<M(y)\). If \(R(x)=R(y)\) and \(\lambda(x)<\lambda(y)\), either \(M(x)<M(y)\) or \(M(x)=M(y)\), in which case the rank inequality decides. Both cases prove the assertion. The condition \(x\in W'\) is essential to the equality \(R(x)=M(x)\); we do not assert order preservation for arbitrary exterior pairs.
\by{the index inequalities from Step 1; rank comparison}
\end{proof}

\subsection{The two child updates}

For an edge \(W\to W'\) on a same side, with cutoff \(S\), let the \textbf{carry} be empty if \(O\setminus S=\varnothing\). Otherwise let \((d,r)\) be the least old key of a job of \(O\setminus S\), and use (5.1). Append the \textbf{injection} \[
I=D_a\setminus P_b
\tag{5.2}
\] as the last entry, then prune on \(W'\). The injection is an upset: it is the intersection of \(D_a\) and the complement of the downset \(P_b\).

On a switched side, discard the old list. Section 6 will construct an upset \(G\) in the working frame satisfying \[
X=W\setminus D_a\subseteq G,\qquad
G\cap\{x:x\text{ is scheduled at or before }a\}=\varnothing.
\tag{5.3}
\] The cutoff \(S\) is a rank prefix of \(X\). Choose \(h\) so that \(X\setminus S=X\cap\{\lambda\geq h\}\), taking \(h=N+1\) if \(X\subseteq S\). Use the \textbf{switch pair} \[
\bigl(G\cap\{\lambda\geq h\},D_a\bigr)
\tag{5.4}
\] and prune it on \(W'\).

\begin{numbered}{Lemma}{5.3}{qualifiers and child (I1)}

For either update, the union before pruning intersects \(W\) in exactly \(W\setminus S\), and its key order on that set agrees with the side's priority order. After pruning, its qualifiers in \(W'\) are exactly \(W'\setminus S\). With this list, \(W'\) satisfies (I1) in its tag frame.

\end{numbered}
\uses{Lemmas 4.2, 4.3, 4.5, and 5.1; invariant (I1); the stipulated switch-upset property (5.3).}

\begin{proof}
\step{1}{the same update has the stated qualifiers and order}
The two working sides, same and switched, exhaust the construction. On a same side, \(O\) is promoted in key order. By Lemma 5.1, the carry selects exactly \(O\setminus S\) within \(W\) and preserves its old key order. By Lemma 4.2, \[
W\cap I=(W\cap D_a)\setminus P_b=Y.
\] The sets \(Y\) and \(O\) are disjoint, since \(O\subseteq P_b\). Also \(W=S_0\dot\cup O\dot\cup Y\), and \(S\) is \(S_0\) together with the promoted part of \(O\). Thus the new qualifiers in \(W\) are \((O\setminus S)\dot\cup Y=W\setminus S\). Placing the injection last puts \(Y\) after the carry qualifiers, with rank order inside \(Y\), exactly as in the priority order.
\by{Lemmas 4.2, 4.3, and 5.1; invariant (I1)}

\step{2}{the switched update has the stated qualifiers and order}
On a switched side, \(G\supseteq X\) implies that the first entry selects exactly \(X\setminus S\) within \(X\). The second entry selects \(W\cap D_a=W\setminus X\). Its larger index puts every job of this second group after those of \(X\setminus S\), regardless of any overlap with the first entry. Keys use rank within each group, agreeing with the priority order. The union in \(W\) is again \(W\setminus S\).
\by{property (5.3); switch-pair definition (5.4)}

\step{3}{pruning preserves the child qualifiers and proves (I1)}
Pruning changes no qualification within \(W'\). Therefore the child's old high set is \(O'=W'\setminus S\). Equation (4.6) gives \(O'\subseteq P_{b'}\) and \(W'\setminus O'=W'\cap S\subseteq D_{a'}\), exactly child (I1).
\by{Lemma 4.5; the pruning definition}
\end{proof}

\subsection{Runs and a capacity witness}

A \textbf{run} is a chain along a branch with a fixed tag, starting at the root or at a child created by a switched update, and proceeding only through same-side updates. Its initial entries are the root entry or the two entries of its initial switch pair. Every other entry is descended from an injection appended during the run.

\begin{numbered}{Lemma}{5.4}{entry shape}

Each entry is an initial entry of its run or an injection appended in that run, intersected with at most one rank suffix in the run's frame.

\end{numbered}
\uses{definitions of carry, pruning, and a run.}

\begin{proof}
\step{1}{the initial entries satisfy the shape invariant}
Induct on the number of same-side updates since an entry was created in its run. At creation the entry is its initial or injection set, with no additional suffix. Assume that every current entry has the stated form.
\by{the initial-entry and injection definitions}

\step{2}{each update preserves a single effective suffix}
Subsequent operations keep or delete an entry, or intersect it with a rank suffix. Intersecting two suffixes gives the suffix at the maximum of their thresholds: \[
\{\lambda\geq r\}\cap\{\lambda\geq r'\}
=\{\lambda\geq\max(r,r')\}.
\] All thresholds in one run use the same rank. Thus at most one suffix is needed in the resulting expression. A switched update starts a new run and deletes the earlier list.
\by{the induction hypothesis; the maximum-threshold identity}
\end{proof}

\begin{numbered}{Lemma}{5.5}{at most \(m\) injections and \(m+2\) entries}

Every pruned list contains at most \(m\) injections and at most \(m+2\) entries.

\end{numbered}
\uses{Lemmas 4.5, 5.1, 5.3, and 5.4; invariant (I1); injection definition (5.2).}

\begin{proof}
A switched update starts a new run with a list of at most two entries and no injection (Section 5.1). So consider a
same-side update $W\to W'$ with edge data $a',b',S$ as in (4.5), and let $\mathscr I$ be the finite set of injection
entries that survive in the pruned list of $W'$, distinguished by their times of creation.
\step{1}{two entries of $\mathscr I$ are disjoint on $W$} Let $K,K'\in\mathscr I$ with $K$ created before $K'$, and
let $U\to U'$ be the edge of the run at which $K'$ was appended, so $W\subseteq U$. At that time $K$ belonged to the
carry of $U$. The qualifiers of the carry inside $U$ are old qualifiers of $U$, hence lie in $P_{b_U}$. The injection
appended at that edge is $D_{a_U}\setminus P_{b_U}$, which meets $U$ only outside $P_{b_U}$. So $K\cap K'\cap
U=\varnothing$ at that time. Later operations only delete entries or intersect them with rank suffixes, and
$W\subseteq U$. Hence $K\cap K'\cap W=\varnothing$. \by{Lemma 5.1; invariant (I1) at $U$; (5.2); Lemma 5.4}
\step{2}{every $K\in\mathscr I$ contains a job of the slot $b'$} $K$ survived pruning on $W'$, so $K\cap
W'\neq\varnothing$; let $x_K$ be the least-rank job of $K\cap W'$. It is a qualifier in $W'$, hence a high job of the child, so
$x_K\in P_{b'}$: some job of the slot $b'$ is a strict successor of $x_K$. $K$ is an upset of all of $J$, so it
contains that successor. Let $f(K)$ be the job of least rank in $b'$ that is a successor of $x_K$. Then $f(K)\in
K\cap b'$. \by{Lemma 5.3; (4.6); every entry is an upset of $J$}
\step{3}{$f\colon\mathscr I\to b'$ is injective} The slot $b'$ is a slot of $W$. If $f(K)=f(K')$ for $K\neq K'$, this
job lies in $K\cap K'\cap W$, against Step 1. \by{Lemma 4.5; Step 1}
\step{4}{the count} By Step 3, $|\mathscr I|\leq|b'|=m$. Every entry is an initial entry of its run or descends from
an injection appended in the run, and a run has at most two initial entries. So a pruned list has at most $m+2$
entries. \by{Step 3; Lemma 5.4; definition of a run, Section 5.2}

\noindent\emph{Where $m$ enters:} in Step 4, as the size $|b'|=m$ of one slot. The witness $f(K)$ lies outside $W'$,
which is why the entries must be upsets of all of $J$ (Step 2). No statement about formula size or deeper nodes is
used.
\end{proof}

Before pruning, a newly appended list may have \(m+3\) entries. Every carry used in the estimates below comes from a pruned list and contains at most \(m\) injections and at most two initial entries.

\section{Direction changes and bounded formulas}

\begin{numbered}{Lemma}{6.1}{local flattening of the initial pair}

Suppose a run starts on a switched side of a node \(Q_0\). In the run's frame its initial pair is \[
G_0\cap\{\lambda\geq h_0\},\qquad D_{a_0},
\] where \(G_0\) is an upset containing \(Q_0\setminus D_{a_0}\). Let \(U\subseteq Q_0\) be a node in this run. Suppose a same-side carry out of \(U\) still retains the first initial entry, now in the form \(G_0\cap\{\lambda\geq k\}\). On \(U\), the union of the carried descendants of the initial pair agrees with \[
D_{a_0}\cup\{\lambda\geq k\}.
\tag{6.1}
\]

\end{numbered}
\uses{Lemmas 5.1 and 5.4; the pruning and run definitions; switch pair (5.4).}

\begin{proof}
\step{1}{a surviving first identity leaves the second untrimmed or absent on the node}
A carry retains a suffix and trims only its first entry. If it ever starts at the second initial entry or later, the first initial entry is permanently deleted. Therefore, as long as the first survives, the second has never been trimmed. The second either remains exactly \(D_{a_0}\) or was deleted by pruning on an earlier node containing \(U\), in which case it is disjoint from \(U\).
\by{Lemma 5.1; pruning; the nesting of nodes in a run}

\step{2}{the two formulas agree on the whole node}
The disjoint sets \(U\cap D_{a_0}\) and \(U\setminus D_{a_0}\) partition \(U\). Within \(U\cap D_{a_0}\), the second entry, when relevant, selects every job. Within \(U\setminus D_{a_0}\), all jobs belong to \(G_0\) because \(U\subseteq Q_0\); there the first entry selects exactly the rank suffix. This proves equality on \(U\). Both terms in (6.1) are upsets. No equality outside \(U\) is claimed or used.
\by{the two parts $U\cap D_{a_0}$ and $U\setminus D_{a_0}$; $U\subseteq Q_0$}
\end{proof}

\begin{numbered}{Lemma}{6.2}{short extension of an earlier low set}

For an edge \(U\to W\) with low cutoff \(H\subseteq U\) in its working frame, there exists a downset \(H^*\subseteq J\) in that frame such that \[
H^*\cap U=H,
\tag{6.2}
\] and \(H^*\) has a formula of at most \(3m+3\) leaves.

\end{numbered}
\uses{Lemmas 2.6, 5.3, 5.4, 5.5, and 6.1; the edge-cutoff construction.}

\begin{proof}
\step{1}{the switched-edge extension uses two leaves}
An edge is either switched or same; these two cases exhaust the edge construction. Use \(a_U,b_U\) for the boundaries of \(U\) in this edge's working frame.

If the edge is on a switched side, \(H\) is a rank prefix of \(U\setminus D_{a_U}\). Let \(r\) be its largest rank, or \(0\) for an empty prefix. Then \[
H^*=(J\setminus D_{a_U})\cap\{\lambda\leq r\}
\tag{6.3}
\] is a downset described with two leaves and has the required intersection.
\by{the switched cutoff definition; cone and rank closure}

\step{2}{express the same-edge low set by its carry}
For a same-side edge, let \(\Phi\) be the carry union before the injection. By Lemma 5.3, or directly by its same-side calculation, \[
H=(U\cap P_{b_U})\setminus\Phi.
\tag{6.4}
\] The carry contains at most \(m\) injection descendants. Each has at most three leaves: \(D\), the complement of \(P\), and at most one rank suffix.
\by{Lemma 5.3; Lemmas 5.4 and 5.5}

\step{3}{replace the initial contribution by at most two leaves on the node}
For the initial contribution, the cases are a root run, a switched run retaining its first initial entry, a switched run retaining only its second, or no retained initial entry. Its initial contribution can be replaced on \(U\) by an upset using at most two leaves, as follows. In a root run, the root entry with a suffix simplifies to a rank suffix (or a constant). In a switched run, if the first initial descendant survives, replace it and the second if present by (6.1). If the first does not survive, the only possible initial contribution is the second entry with at most one suffix, using two leaves. If no initial entry remains, omit that contribution. These cases exhaust the run's initial entries.
\by{Lemma 6.1; Lemma 5.4; the run definitions}

\step{4}{construct the downset extension and count all leaves}
Thus there is an upset \(\Phi'\) agreeing with \(\Phi\) on \(U\), with at most \[
3m+2
\tag{6.5}
\] leaves. If the carry is empty, the constant false is an allowed one-leaf substitute and also satisfies this bound. Set \[
H^*=P_{b_U}\cap(J\setminus\Phi').
\tag{6.6}
\] The complement of an upset is a downset, so this is a downset. Equations (6.4) and the agreement on \(U\) prove (6.2). Its leaf count is at most \(1+(3m+2)=3m+3\). The displayed extension uses the replacement initial contribution from (6.1), with the stated leaf count independent of earlier depth.
\by{Steps 2--3; Lemma 2.6; (6.4)}
\end{proof}

\begin{numbered}{Lemma}{6.3}{arrangement at a switch}

Suppose a switched side of \(W\) produces a nonempty child. Then \(W\) is not the root. Let \(U\) be its parent. In the current, switched frame write \(a_-,b_-\) for the boundaries of \(U\), \(v_-\) for its center, and \(a,b\) for the boundaries of \(W\). Then \(a\) is a genuine slot and \[
a_-<v_-\leq a<b\leq b_-.
\tag{6.7}
\]

\end{numbered}
\uses{Lemma 4.5; the frame-reversal map.}

\begin{proof}
\step{1}{the node has a parent and the old edge lies before its center}
The root has no switched-side child by Lemma 4.5. The edge \(U\to W\) was taken in the tag frame of \(W\), opposite to the current frame. In that old frame, \(W\) is a gap before \(U\)'s center. Its old right boundary is a slot and is at or before that center; its old left boundary is at or after \(U\)'s left boundary.
\by{Lemma 4.5; the child construction}

\step{2}{reversal gives the displayed order of boundaries and center}
Reversal sends its old right boundary to \(a\), its old left boundary to \(b\), the old right boundary of \(U\) to \(a_-\), and the old left boundary to \(b_-\). These inequalities give (6.7).
\by{the decreasing bijection $t\mapsto T+1-t$}
\end{proof}

Let \(H\) be the old cutoff of edge \(U\to W\) and let \(H^*\) be its extension from Lemma 6.2. It is a downset in that old frame, hence an upset in the current one. Define \[
G
=
D_{a_-}
\cap(J\setminus\widehat P_a)
\cap(J\setminus\widehat P_{v_-})
\cap H^*.
\tag{6.8}
\]

\begin{numbered}{Lemma}{6.4}{the switch upset}

The set (6.8) is an upset in the current frame, contains all of \(X=W\setminus D_a\), and contains no job scheduled at or before \(a\).

\end{numbered}
\uses{Lemmas 2.5, 4.3, 4.4, 4.5, 6.2, and 6.3; definition (6.8).}

\begin{proof}
\step{1}{the switch set is a global upset}
Each factor of (6.8) is an upset: a successor cone, two complements of downsets, and \(H^*\).
\by{Lemma 6.2; cone closure; frame reversal}

\step{2}{every job of $X$ satisfies all four factors}
Let \(x\in X\). If \(x\) had been high at the edge \(U\to W\), (4.6) would put it in the old predecessor cone of that edge's right boundary. That boundary is current \(a\), and its old predecessor cone is current \(D_a\), contradicting \(x\in X\). Thus \(x\in H\). Every old edge cutoff lies in the old \(P\)-cone of \(U\)'s old right boundary, by Lemma 4.3. This cone is current \(D_{a_-}\), so \(x\in D_{a_-}\). It is after \(a\) and hence after \(v_-\); Lemma 2.5 excludes it from both weak predecessor cones in (6.8). Finally \(x\in U\) and (6.2) give \(x\in H^*\). This verifies every factor, proving \(X\subseteq G\).
\by{Lemmas 4.3, 4.5, 2.5, and 6.2}

\step{3}{a putative early qualifier is localized to the old node}
Now suppose \(z\in G\) is at or before \(a\). Its membership in \(D_{a_-}\) puts it strictly after \(a_-\). From (6.7) it is also before \(b_-\). Therefore \(z\in U\), and only after this localization may we use (6.2), obtaining \(z\in H\). The two possible time ranges are at or before \(v_-\), and strictly after \(v_-\) through \(a\).
\by{Lemmas 6.3 and 6.2}

\step{4}{the interval through the old center contains no such qualifier}
If \(z\) is at or before \(v_-\) in the current frame, it is at or after the center in the old frame. At the center it belongs to \(\widehat P_{v_-}\). Otherwise it is a low job after the old center: \(H\) is contained in that side's center cutoff, and the center separates that cutoff on \(U\). Thus it lies in old \(D_{v_-}\), which is current \(P_{v_-}\). Either possibility contradicts its exclusion from \(\widehat P_{v_-}\).
\by{the old center-separator property; Step 3}

\step{5}{the interval from the old center to the current boundary contains no such qualifier}
If \(z\) is strictly after \(v_-\) and at or before \(a\), in the old frame it is on the side before the center, at or after the old edge's right boundary \(a\). At that boundary it belongs to its weak cone; after that boundary it is low for \(H\) and belongs to old \(D_a\) by Lemma 4.4. Hence in either event it is in current \(\widehat P_a\), another contradiction.
\by{Lemma 4.4; the weak-cone definition; Step 3}

\step{6}{the needed exterior separator facts persist in the current witness}
These separator statements were established in the witness of \(U\). Subsequent changes on this branch before constructing (6.8) occurred only inside \(W\), strictly between \(a,b\). Jobs at or before \(a\), the boundary slots and center used here, and the fixed sets \(H,H^*\) have not changed. Thus the separator information used for \(z\) remains valid in the current witness. This proves all three assertions.
\by{the nesting of replacements strictly inside $W$}
\end{proof}

This supplies exactly (5.3), for the whole parent's \(X\), not merely its intersection with a chosen child. It also proves that the switch pair has no qualifier at or before \(a\): the first entry is contained in \(G\), and the second in \(D_a\).

\begin{numbered}{Lemma}{6.5}{individual and aggregate leaf bounds}

Every list entry is described by at most \(3m+7\) leaves. More particularly:

\begin{itemize}
\tightlist
\item
  an injection descendant has at most three leaves;
\item
  the root's initial descendant has at most two leaves, and the second initial descendant of a switched run at most two;
\item
  only the first initial descendant of a switched run may need \(3m+7\) leaves.
\end{itemize}

Every pruned-list union and every carry union has a description with at most \[
A_m=6m+9
\tag{6.9}
\] leaves.

\end{numbered}
\uses{Lemmas 2.6, 5.4, 5.5, and 6.2; definitions (5.2), (5.4), and (6.8).}

\begin{proof}
\step{1}{every entry has its stated leaf budget}
The first two bounds follow from Lemma 5.4 and their displayed definitions. Formula (6.8) has three cone factors in addition to \(H^*\), so it has at most \(3+(3m+3)=3m+6\) leaves. A first switch entry adds one rank suffix, giving \(3m+7\). Further suffix intersections combine with that same suffix; they do not add indefinitely. A new run discards the old list.
\by{Lemmas 5.4 and 6.2; definitions (5.2), (5.4), and (6.8)}

\step{2}{every list and carry union uses at most $6m+9$ leaves}
For an aggregate union in a switched run, bound its at most \(m\) injections by \(3m\), its possible first initial entry by \(3m+7\), and its possible second by \(2\). Their sum is \(6m+9\). If some entries are absent, omitting them only reduces the bound. For a root run, the bound is smaller. An empty union is the constant false and costs one leaf. The same reasoning applies to a carry from a pruned list, because it cannot increase the number of entries and only combines a suffix into its first entry.
\by{Lemma 5.5; Step 1; the one-leaf empty set}
\end{proof}

This aggregate count is why the final leaf budget is linear in \(m\). Multiplying the uniform per-entry bound by \(m+2\) would give a valid but unnecessarily quadratic estimate.

\section{Preserving the exterior boundary invariants}

The preparatory rescheduling of a node preserves its existing invariants. In (I1) its job set and all cone/list sets are fixed. In (I2) jobs at or before its left boundary and the boundary slot do not move. In (I3) the quantifier already includes every full replacement schedule of its interval. Thus choosing a center and normalizing the two sides does not invalidate the induction hypotheses.

Fix an edge \(W\to W'\) with working-frame boundaries \(a,b\), normalized side \(E\), and data \(a'=s_i\), \(b'=s_{i+1}\), \(S=S_i\). The child's tag is this working frame.

\begin{numbered}{Lemma}{7.1}{child (I2)}

Every qualifier of the child's pruned list scheduled at or before \(a'\) belongs to \(\widehat P_{a'}\).

\end{numbered}
\uses{Lemmas 2.5, 3.4, 3.6, 4.3, 4.4, 5.1, 5.3, and 6.4; invariants (I1), (I2), and (I3).}

\begin{proof}
\step{1}{reduce to the pre-pruning qualifiers and split their possible times}
Pruning can only remove qualifiers, so take a qualifier \(z\) of the list before pruning at or before \(a'\). Its time is either after \(a\) or at or before \(a\); these cases are exhaustive.
\by{the pruning definition}

\step{2}{a qualifier after the old boundary satisfies child (I2)}
First suppose that \(z\) is after \(a\). The interval between \(a\) and \(a'\) lies in \(E\), since \(a'\) is before the center. Thus \(z\in E\subseteq W\) and it is high for \(S_i\) by Lemma 5.3. If it occupies \(a'\), it belongs to \(\widehat P_{a'}\). Otherwise it is a high job before the separator \(a'\) and therefore belongs to \(P_{a'}\).
\by{Lemma 5.3; Lemma 4.4}

\step{3}{an earlier qualifier is carried from the parent boundary}
Now suppose that \(z\) is at or before \(a\). No injection qualifier can be here, since the injection is contained in \(D_a\). Neither switch-pair entry can have a qualifier here, by Lemma 6.4 and the last paragraph following its proof. Therefore this is a same update, and \(z\) qualifies through its carry. Lemma 5.1 implies that it was an old qualifier. Parent (I2) gives a job \(y\) in slot \(a\) with \(z=y\) or \(z\prec y\). The carry union is an upset, so \(y\) also qualifies through the carry. In particular \(a\) is genuine.
\by{Lemmas 2.5, 6.4, and 5.1; parent (I2)}

\step{4}{the unchanged boundary is settled; later positions retain a high carry}
If \(a'=a\), the old (I2) assertion is already the desired one. Suppose \(a'>a\). Follow every actual advance of the walk between \(a\) and \(s_i=a'\), including positions within a promotion that were not marked. Such an advance uses a cutoff \(S_j\subseteq S_i\). By Lemma 5.3, the carry at the final edge selects in \(W\) exactly \(O\setminus S_i\), so every carry qualifier in \(E\) is high for every cutoff encountered in this part of the walk.
\by{Lemma 5.3; the nested side cutoffs}

\step{5}{fix the first actual boundary advance}
Let the first advance leave \(a\) for the slot \(t\) of an exception \(x\), during promotion \(j\leq i\). We verify the required boundary exchange carefully.
\by{the side-walk construction}

\step{6}{the exception has every predecessor before the old boundary}
Every predecessor of \(x\) in \(J\) is strictly before \(a\). Every predecessor is before \(t\) by feasibility. A predecessor in slot \(a\) would imply \(x\in D_a\). A predecessor strictly between \(a\) and \(t\) would be a job of \(E\) in \(S_j\) (because \(S_j\) is a priority prefix and \(x\in S_j\)). Since \(t\) is the earliest exception slot, that predecessor would lie in \(D_a\), again implying \(x\in D_a\). Both contradict that \(x\) is an exception. This checks the assertion globally, including exterior jobs.
\by{priority-prefix closure; earliest-exception choice; transitivity}

\step{7}{the exception has smaller old key than the boundary qualifier}
The old key of \(x\) is smaller than that of \(y\). Since \(x\notin D_a\), (I1) gives \(x\in O\). It does not belong to \(S_0\subseteq D_a\), so its presence in \(S_j\subseteq S_i\) means it was promoted from \(O\). The carry is nonempty in this case; let \((d,r)\) be its threshold, the least key in \(O\setminus S_i\). The promotion order gives \(\operatorname{key}(x)<(d,r)\), while Lemma 5.1 gives \(\operatorname{key}(y)\geq(d,r)\).
\by{invariant (I1); Lemmas 4.3 and 5.1}

\step{8}{parent (I3) forbids their exchange}
The exchange of \(x\) with \(y\) is forbidden by parent (I3). The current schedule in \(W\) is a full replacement of its old witness, \(x\in W\setminus D_a\), and \(y\) is an old qualifier at \(a\) with the requisite key inequality.
\by{parent (I3); Step 7}

\step{9}{a successor of the boundary qualifier must occur by the target}
Consequently \(y\) has a successor at or before \(t\). If it did not, swapping \(x,y\) would be feasible: the global predecessors of \(x\) are before \(a\) by Step 6, and its successors are after \(t\); the predecessors of \(y\) are before \(a\) by the old schedule, and its successors would be after \(t\) by the supposition. These facts also exclude comparabilities between a moved job and any companion in either destination slot. All other relations are unchanged. This would contradict Step 8.
\by{Steps 6 and 8; the full-schedule precedence conditions}

\step{10}{the carry has a successor at the first target}
Such a successor \(y_1\) of \(y\) is after \(a\) and at or before \(t\), hence in \(E\), and qualifies through the upset carry. If it is at \(t\), the first advance has propagated a carry qualifier from \(a\) to \(t\). If it is strictly before \(t\), apply Lemma 3.4(2) with \(M=L=S_j\cap E\), starting at boundary \(a\); the order and cutoff restrictions are as in the paragraph after Lemma 3.6. Its hypotheses hold: condition (a) at the exterior boundary is empty, \(x\) is the earliest exception and is low, and \(y_1\) is high. The lemma gives a successor of \(y_1\) at \(t\). The upset carry contains it. Thus in all cases \(y\) has a carry-qualifying successor at the first target. This target job cannot be \(x\), which is low.
\by{upset closure; Lemma 3.4(2) on the restricted domain}

\step{11}{propagation through every later advance proves child (I2)}
For a later advance from a genuine slot \(s\) to \(t'\), its cutoff is some \(S_j\subseteq S_i\). A carry qualifier at \(s\) is high for that cutoff. Lemma 3.4(2), with \(M=L=S_j\cap E\), gives a strict successor at \(t'\); it still qualifies through the carry because the union is an upset. Following all actual advances gives a job in \(a'\) strictly above \(y\). Since \(z=y\) or \(z\prec y\), it follows that \(z\in P_{a'}\).

This proves (I2) in both cases. The sole use of the external exchange invariant was at the first advance; all later exchanges are internal to normalized \(E\).
\by{Lemma 3.4(2); upset closure; transitivity}
\end{proof}

\begin{numbered}{Lemma}{7.2}{child (I3)}

Replace the interval of \(W'\) by any full schedule of \(W'\). Let \(x\in W'\setminus D_{a'}\) and let \(y\) be a qualifier of its pruned list in slot \(a'\), with child's key of \(x\) smaller than that of \(y\). Exchanging \(x,y\) cannot give a full schedule of \(J\).

\end{numbered}
\uses{Lemmas 2.4, 2.5, 3.3, 5.1, 5.2, 5.3, and 6.4; invariants (I1) and (I3).}

\begin{proof}
\step{1}{the child key comparison holds before pruning}
Child (I1), already proved in Lemma 5.3, makes \(x\) a qualifier. Lemma 5.2 shows that the inequality between its key and that of \(y\) holds already before pruning. The construction gives \(a'\geq a\); the two cases are \(a'>a\) and \(a'=a\).
\by{Lemmas 5.3 and 5.2}

\step{2}{an interior left boundary uses normalized-side minimality}
If \(a'>a\), the slot \(a'\) belongs to \(E\), so both \(x\) and \(y\) belong to \(W\). By Lemma 5.3 they are high for \(S_i\), and their pre-pruning key order agrees with the side's priority order. Thus \(x\) has better priority than \(y\). The arbitrary replacement of \(W'\) changes only slots of \(E\) strictly after \(a'\) and, by Lemma 2.4, gives a full schedule of \(E\). If the exchange were feasible in \(J\), its restriction would be a full schedule of \(E\) violating Lemma 3.3 at \(r=a'\). This is impossible.
\by{Lemmas 5.3, 2.4, and 3.3}

\step{3}{an inherited left boundary uses the parent exchange invariant}
If \(a'=a\), a switched update has no qualifier at \(a\), so that case has no \(y\). In a same update, \(x\notin D_a\) by hypothesis and \(y\notin D_a\) by Lemma 2.5. Neither can qualify through the injection. They both qualify through the carry, and Lemma 5.1 transfers the strict key inequality to the old list. Also \(x\in W\setminus D_a\), while \(y\) is an old qualifier in \(a\). The preparatory changes in \(W\) followed by the chosen schedule of \(W'\) form a full schedule of \(W\) in its interval. Parent (I3) therefore forbids exactly this exchange.
\by{Lemmas 6.4, 2.5, and 5.1; parent (I3)}
\end{proof}

\begin{numbered}{Proposition}{7.3}{valid construction, with no circular induction}

Starting from any full schedule of \(J\), every node and update described in Sections 4--7 is well-defined. Each pruned list has at most \(m+2\) entries with the bounds of Lemma 6.5, each node satisfies (I1)--(I3), and every branch terminates.

\end{numbered}
\uses{Lemma 2.4; Corollary 3.7; Lemmas 4.3--4.5, 5.3, 5.5, 6.1--6.5, 7.1, and 7.2; the root invariants.}

\begin{proof}
\step{1}{construct the center, walks, and edge data by depth induction}
Induct from the root on node depth. Its data and invariants were verified in Section 4. Assume the data of a node and all its ancestors have been constructed and validated. Corollary 3.7 and Lemma 2.4 supply its center and normalized sides. Lemmas 4.3--4.5 give the walks, marks and child edge data.
\by{the root invariants; Corollary 3.7; Lemmas 2.4 and 4.3--4.5}

\step{2}{all list data used to form a child already exist}
A same update uses only this node's existing list. A switched update uses the earlier edge from its parent: its cutoff and, for a same edge, its already formed carry. Lemma 6.1 uses that earlier run's initial entries, which likewise already exist. Lemma 6.2 gives the earlier low-set extension, and Lemmas 6.3--6.4 give the current \(G\). None asks for a descendant's invariant or formula.
\by{Lemmas 6.1--6.4; the ancestor induction hypothesis}

\step{3}{the child satisfies all list and boundary assertions}
The new list is an upset list; Lemma 5.3 proves its qualifier identity and (I1). Lemma 5.5 proves its length using only that identity, global upsets, and existing edge labels. The leaf estimate of Lemma 6.5 uses the length bound and flattening, not an unproved boundary invariant. Finally Lemmas 7.1--7.2 establish (I2) and (I3) from the parent's corresponding hypotheses. This closes the induction.
\by{Lemmas 5.3, 5.5, 6.5, 7.1, and 7.2; parent invariants}

\step{4}{the hierarchy construction terminates}
Every child omits the center of its parent and has a full schedule; its number of slots is strictly smaller. Thus no branch can have more than \(N/m\) nonempty nodes. Each node has finitely many children, so the resulting tree is finite.
\by{strict decrease of the integer slot count; finite branching}
\end{proof}

\section{Global descriptions and the structure theorem}

Fix a node \(W\) and a side, with working-frame boundaries \(a,b\), center \(v\), open part \(E\) and cutoffs \(S_i\).

For a same side, let \(\Phi_i\) be the carry union determined by \(O\setminus S_i\), before adding an injection, and put \[
\widetilde S_i=P_b\setminus\Phi_i.
\tag{8.1}
\] For a switched side, let \(r_i\) be the last promoted rank in \(X\), or \(0\) before any promotion, and put \[
\widetilde S_i=(J\setminus D_a)\cap\{\lambda\leq r_i\}.
\tag{8.2}
\] These predicates are defined for all cutoffs, including those associated with empty gaps.

\begin{numbered}{Lemma}{8.1}{cutoff predicates}

For every side and every \(i\), \[
\widetilde S_i\cap W=S_i.
\tag{8.3}
\] The predicate has at most \(A_m+1\) leaves, where \(A_m=6m+9\).

\end{numbered}
\uses{Lemmas 2.6, 4.3, 5.1, and 6.5; definitions (8.1) and (8.2).}

\begin{proof}
\step{1}{the same-side predicate agrees with its cutoff and meets the budget}
The two working sides, same and switched, exhaust the construction. On a same side, the carry selects \(O\setminus S_i\) in \(W\) by Lemma 5.1. Since \(W\cap P_b=S_0\dot\cup O\) and \(S_i\) consists of \(S_0\) and the promoted part of \(O\), subtracting that carry gives exactly \(S_i\). Its union has at most \(A_m\) leaves by Lemma 6.5, so adding the \(P_b\) factor gives \(A_m+1\).
\by{Lemmas 4.3, 5.1, 6.5, and 2.6}

\step{2}{the switched-side predicate agrees with its cutoff and uses two leaves}
On a switched side (8.2) is exactly the promoted rank prefix of \(X\) within \(W\), and uses two leaves. This is below \(A_m+1\).
\by{the rank-prefix definition; (8.2)}
\end{proof}

The equalities (8.3) specify these predicates on \(W\); subsequent uses first restrict membership to \(W\).

\begin{numbered}{Lemma}{8.2}{global description of a child}

For a child \(W'\) with edge data \(a',b',S=S_i\) and pruned list \(\mathcal K'\), one has the following equalities of subsets of \textbf{all} \(J\): \[
W'\setminus S
=
\mathcal Q(\mathcal K')
\cap P_{b'}
\cap(J\setminus\widehat P_{a'}),
\tag{8.4}
\] and \[
W'\cap S
=
D_{a'}
\cap P_b
\cap(J\setminus\widehat D_{b'})
\cap(J\setminus\widehat D_v)
\cap\widetilde S_i.
\tag{8.5}
\]

\end{numbered}
\uses{Lemmas 2.5, 4.3, 4.4, 4.5, 5.3, 7.1, and 8.1; the center-separator property.}

\begin{proof}
\step{1}{the high-part identity holds globally}
For (8.4), a high job of \(W'\) qualifies by Lemma 5.3, belongs to \(P_{b'}\) by (4.6), and is after \(a'\), excluding membership in \(\widehat P_{a'}\). Conversely, let a job satisfy the right-hand side. Its \(P_{b'}\) membership puts it before \(b'\). If it were at or before \(a'\), child (I2) would put it in \(\widehat P_{a'}\), contrary to the last factor. Hence it is strictly between \(a',b'\), so in \(W'\). The qualifier identity then makes it high.
\by{Lemmas 5.3, 4.5, 2.5, and 7.1}

\step{2}{every selected low-part job lies in the child}
For (8.5), first let \(z\) satisfy the right-hand side. Membership in \(D_{a'}\) puts it after \(a'\geq a\), and membership in \(P_b\) puts it before \(b\). Thus \(z\in W\). Now, and only now, (8.3) shows \(z\in S_i\). Since \(S_i\subseteq S_*\), a job in the center would belong to \(\widehat D_v\) and a job after the center would belong to \(D_v\) by the center's separator condition. Both are excluded. Therefore \(z\) belongs to the side \(E\). If it were in \(b'\) or after it on that side, it would belong to \(\widehat D_{b'}\) by its being low and by Lemma 4.4. For \(b'=v\), the preceding center exclusion already handles this case. Thus \(z\) lies strictly between \(a',b'\), and belongs to \(W'\cap S\).
\by{Lemmas 8.1, 4.3, and 4.4; the center-separator property}

\step{3}{every low child job satisfies the low-part formula}
Conversely, a low job of \(W'\) is in \(D_{a'}\cap P_b\) by (4.6), is in \(\widetilde S_i\) by (8.3), and occurs before \(b'\) and \(v\). Lemma 2.5 excludes it from both indicated weak successor cones. It satisfies every factor of the right-hand side.
\by{Lemmas 4.5, 8.1, and 2.5}
\end{proof}

The proof of (8.5) first derives membership in the parent from its cone factors, which permits the use of (8.3). Its formula therefore uses exactly the five displayed factors.

\begin{numbered}{Lemma}{8.3}{relative descriptions at every mark}

For each mark \(z\) of \(W\), its original-time prefix inside \(W\) has a relative description.

\end{numbered}
\uses{Lemmas 2.6, 3.1, 4.4, and 8.1; the center-cutoff definition.}

\begin{proof}
\step{1}{the center formula selects its open side}
In the working frame of a side containing \(z\), or either side for the center, put \[
H_*=
\begin{cases}
P_b,&\text{same side},\\
J\setminus D_a,&\text{switched side}.
\end{cases}
\] Then \(H_*\cap W=S_*\). By Lemma 3.1 the formula \[
F_v=B_v(H_*)
\tag{8.6}
\] selects inside \(W\) exactly the open side \(E\) before the center.
\by{the center-cutoff definition; Lemma 3.1}

\step{2}{the side-mark formula selects its working-frame prefix}
For a side mark \(z=s_i\) with \(i\geq1\), Lemma 4.4 says \(z\) separates \(S_i\) on \(E\). Set \[
F_z=B_v(H_*)\cap B_z(\widetilde S_i).
\tag{8.7}
\] Within \(W\) the first factor selects \(E\). Within \(E\), (8.3) and Lemma 3.1 show that the second selects its jobs before \(z\). These are all the jobs of \(W\) before \(z\) in the working frame.
\by{Lemmas 4.4, 8.1, and 3.1; Step 1}

\step{3}{the stated formula gives the original-time prefix in either frame}
The working frame is either the forward frame or the reverse frame. If that is the forward frame, use \(F_z\). If it is the reverse frame, the original-time prefix is described relatively by \[
(J\setminus F_z)\setminus z.
\tag{8.8}
\] Indeed the jobs of \(W\) partition into its working-frame prefix, the slot \(z\), and its working-frame suffix; the latter is the original-time prefix. Lemma 2.6 permits complementation without increasing the leaf count, and membership in \(z\) is one atom. The same conversion works for the center formula (8.6).
\by{the prefix/slot/suffix partition; Lemma 2.6; time reversal}
\end{proof}

\begin{numbered}{Lemma}{8.4}{uniform leaf counts}

Every child has a global description with at most \(12m+25\) leaves. Every marked prefix has a relative description with at most \(6m+16\) leaves.

\end{numbered}
\uses{Lemmas 2.6, 6.5, 8.1, 8.2, and 8.3; the atomic formula alphabet.}

\begin{proof}
\step{1}{the global child description uses at most $12m+25$ leaves}
Use \(A_m=6m+9\). The qualifier union in (8.4) has at most \(A_m\) leaves, so that equation's right-hand side has at most \(A_m+2\). The right-hand side of (8.5) adds four one-leaf cone tests to a cutoff predicate of at most \(A_m+1\) leaves, giving at most \(A_m+5\). Their union describes \(W'\) globally and costs at most \[
(A_m+2)+(A_m+5)=2A_m+7=12m+25.
\tag{8.9}
\]
\by{Lemmas 6.5, 8.1, and 8.2; Lemma 2.6}

\step{2}{the relative prefix uses at most $6m+16$ leaves in the permitted alphabet}
The center formula \(B_v(H_*)\) uses three leaves. The second factor in (8.7) uses at most \(1+(A_m+1)+1=A_m+3\). Thus (8.7) uses at most \(A_m+6\), and reversal (8.8) adds one leaf. This gives \[
A_m+7=6m+16.
\tag{8.10}
\] Every cone used is attached to a genuine \(m\)-element slot or is a constant for a virtual boundary. Every threshold is in one of the two permitted frames. These are formulas in exactly the alphabet of Section 2, rather than a more expressive language.
\by{Lemmas 8.1 and 8.3; Lemma 2.6; the Section 2 atom definitions}
\end{proof}

\begin{numbered}{Theorem}{8.5}{the \(m\)-machine structure theorem}

For \(m\geq3\), every full schedule of \(J\) on \(m\) machines admits a hierarchy in the sense of Definition 2.9 with \[
K_m=12m+25.
\tag{8.11}
\] Its root has the one-leaf description \(J\), every other node has a global description of at most \(K_m\) leaves, and every marked original-time prefix has a relative description of at most \(6m+16\leq K_m\) leaves. These bounds are independent of the depth.

\end{numbered}
\uses{Definition 2.9; Proposition 7.3; Lemmas 8.2, 8.3, and 8.4.}

\begin{proof}
\step{1}{the construction gives the required node and gap structure}
Apply the construction validated in Proposition 7.3. It marks at least the center of every nonempty node, and its children are exactly the nonempty gaps.
\by{Proposition 7.3; Definition 2.9}

\step{2}{the descriptions and finiteness verify every hierarchy clause}
Lemma 8.2 supplies their global descriptions without intersecting with any parent's description. Lemma 8.3 supplies all marked relative prefixes. Lemma 8.4 gives the bounds and Proposition 7.3 gives finiteness. These verify every item of Definition 2.9.
\by{Lemmas 8.2--8.4; Proposition 7.3; Definition 2.9}
\end{proof}


\section{Correctness and running time}

\begin{numbered}{Theorem}{9.1}{correctness of the actual finite algorithm}
For $m\geq3$, run the dynamic program of Section 2 with $m$-element slots and $K=K_m$.
It accepts the padded set $J$ exactly when $J$ has a full schedule, and
returns a full schedule whenever it accepts. Consequently it decides the
original deadline and constructs an optimum by testing deadlines.
\end{numbered}
\uses{Lemmas 2.1, 2.2, 2.3 and 2.10; Proposition 2.8; Theorem 8.5.}
\begin{proof}
\step{1}{the finite search decides the padded instance [Proposition 2.8; Lemmas 2.3 and 2.10; Theorem 8.5]}
Proposition 2.8 turns every accepting state path into a full schedule.
Theorem 8.5 supplies a hierarchy for every full-schedulable $J$, and
Lemma 2.10 maps its marks and child gaps to an accepting state path of the
same finite search. Lemma 2.3 proves both directions of padding and its
removal.

\step{2}{elementary deadlines are decided without large padding [Lemmas 2.1 and 2.3]}
For a nonempty instance reject $T<1$. For $T\geq n$, the topological
order of Lemma 2.1 schedules one job per slot and avoids padding to a
possibly large $mT$. In the remaining branch reject $n>mT$; otherwise
run the padded algorithm.

\step{3}{the deadline trials return an optimum [Steps 1--2; Lemma 2.2]}
Test $T=1,\ldots,n$ in increasing order to optimize. Deadline $n$ is feasible, and the first accepted deadline is the
minimum. Lemma 2.2 also covers the original convention allowing real
start times. Assign the at most $m$ original jobs in a slot to distinct
machines in increasing identifier order.
\end{proof}

The normalized schedules in the structural proof certify completeness;
they need not be computed by the finite algorithm.

\paragraph{Encoding and the start of the padded bound.}
An explicit input uses canonical self-delimiting binary integers for
$m,n$, a signed integer $T$, an explicit list of vertices $1,\ldots,n$,
and a list of distinct arcs with endpoints in that list. Its bit length
$L$ includes delimiters and length fields, so $n\leq L$.
The graph is a DAG; malformed inputs can also be rejected by the
preprocessing below. Arbitrary finite binary vertex names can first be
renamed by repeated input scans. Input reading and that renaming are
charged in $L$. The bound in $N$ begins with the parsed, canonically
numbered padded DAG on $J=[N]$, where $N=mT\geq m$.

\begin{numbered}{Proposition}{9.2}{explicit multitape time bound}
For $m\geq3$ and $K\geq1$, the padded dynamic program, including its
witness on every accepted instance, runs on one deterministic multitape
Turing machine in
\[
 A(m,K)N^{5mK+4m+1}\tag{9.1}
\]
steps, where one effective coefficient bound is
\[
 A(m,K)=2^{40}(m+1)^4(K+1)^4(K512^K)^5.
\]
The machine has a fixed finite number of tapes, independent of $m$ and
$K$.
\end{numbered}
\uses{Lemmas 2.1 and 2.7; Proposition 2.8; the state and transition definitions in Section 2.}
\begin{proof}
\step{1}{at most $B$ formulas and $\mathcal{S}$ states suffice [Lemma 2.7; state definitions]}
Put
\[
 C=K512^K,\quad B=CN^{mK},\quad \mathcal{S}=BN^m,
 \quad h=\lceil\log_2(N+1)\rceil,\quad q=64(m+1)(K+1).
 \tag{9.2}
\]
These are bounds; loops visit only actual records. Lemma 2.7 bounds the
formula count by $B$. Explicitly, the leaf alphabet has at most
$4+8(N+2)+20\binom Nm\leq64N^m$ symbols. For $\ell$ leaves there
are at most $4^\ell$ tree shapes and $2^{\ell-1}$ internal Boolean
labellings. Summing their product with $(64N^m)^\ell$ over
$1\leq\ell\leq K$ gives $B$.
For fixed $W$, a described set $F$ and an $m$-subset $Z$ determine
$L=W\cap F$ and $R=W\setminus(L\cup Z)$. Reject invalid candidates.
Thus at most $B\binom Nm\leq \mathcal{S}$ candidate states remain, including
duplicates.

\step{2}{all records have explicit scan bounds [encoding; elementary tape operations]}
Use a fixed collection of tapes (32 suffice) for the input, matrices,
rank map, formula stack, files, copied records, counters and output.
All variable-length stacks are delimited records on these tapes.
Store the reachability matrix in row-major order. A unary distance
scratch word permits a complete seek to a matrix location in $O(N^2)$
steps, with the scratch and matrix heads advancing together; no binary
counter is incremented for each matrix bit.

Store formula and described-set records in files with record length at
most $qN$. A set record contains its $N$ bits, cardinality, ordinal and
acceptance bit. A current-state record contains its three $N$-bit parts,
ordinal, reached/initial/final flags and predecessor ordinal, also in
at most $qN$ bits. Indeed,
\[
 \lceil\log_2(\mathcal{S}+1)\rceil
 \leq 2+10K+(mK+m)h\leq qh/8\leq qN/8.
\]
Formula codes use $O(K(m+1)h)$ bits, within the same record envelope.
Fetch or update an indexed state by scanning its file and comparing
ordinal fields. This costs $O(q\mathcal{S}N)$ including rewinds and all index
bits. Acceptance queries compare their $N$-bit query against the set
file. Keep saved schedules in a \emph{separate} file, with at most $N$
identifiers per schedule and at most $qN^2$ bits per record. A combined
acceptance/witness lookup costs $O(qBN^2)$.

For explicit slack let $D=2^{20}q^2$. In the scan implementation just
described, each row below costs at most $D$ times its displayed
expression. Scan, copy, compare and rewind use a bounded number of tape
moves per symbol; the records include their counters. At most 64
primitive passes per operation, including the postorder stack scans,
fit this coefficient. No record access has unit cost.

\begin{center}\small
\begin{tabularx}{\linewidth}{@{}p{0.50\linewidth}X@{}}
\toprule
Subroutine, over the whole padded call & Step bound divided by $D$\\
\midrule
Adjacency, reachability and rank & $N^5$\\
Generate and evaluate formulas & $BN^3$\\
Pairwise deduplication & $B^3N$\\
Cardinalities and their ordering & $B^3N^2$\\
Enumerate and initialize states & $B\mathcal{S}(BN^2+N^4)$\\
State-graph propagation & $3B\mathcal{S}^4N$\\
Reconstruct all saved witnesses & $B\mathcal{S}N^2+B^2N^3$\\
Write the selected witness & $N^3$\\
\bottomrule
\end{tabularx}
\end{center}

\step{3}{graph preprocessing costs at most $DN^5$ [Lemma 2.1; Step 2]}
At most $N^2$ arc insertions cost $O(N^4)$ by full matrix seeks.
Floyd--Warshall uses $N^3$ triples, each with a bounded number of
$O(N^2)$ seeks, hence $O(N^5)$ steps. Its invariant is reachability
through the already processed internal vertices. Repeated scans select
the least active vertex without an active predecessor, append it to the
rank order and remove it; all removals fit $O(N^5)$ as well. Lemma 2.1
justifies that rank order.

\step{4}{formula generation and evaluation cost at most $DBN^3$ [Lemma 2.7; Steps 1--2]}
Generate valid Dyck words of length $2(\ell-1)$ by depth-first
extension through valid prefixes; they encode the ordered full binary
trees with $\ell$ leaves. There are $O(K)$ prefix edges per output and
$O(K)$ steps per update, giving $O(K^2)$ delay. Enumerate their Boolean
labellings and leaf parameters by finite counters. A cone leaf needs at most $mN$
reachability queries, each with an $O(N^2)$ seek. At most $K$ leaves
therefore cost $O(mKN^3)$ per formula. Threshold, slot and complement
leaves cost less. Postorder evaluation combines $N$-bit vectors; even
scanning the entire $O(KN)$-bit stack at each node costs $O(K^2N)$.
These parameter factors fit $D$, giving $DBN^3$.

\step{5}{deduplication and ordering fit their table bounds [Step 2]}
For each of at most $B^2$ record pairs, full file seeks fetch and compare
the vectors and mark the later duplicate. Each pair costs $O(qBN)$,
giving $DB^3N$; compact the retained records afterwards. Count each
vector's one-bits and store its cardinality. Repeated selection scans,
even with full record retrieval for every compared pair, cost at most
$DB^3N^2$. Break equal-cardinality ties by ordinal.

\step{6}{state enumeration costs at most $DB\mathcal{S}(BN^2+N^4)$ [state definitions; Steps 1--2]}
For each of $B$ sets and at most $\mathcal{S}$ candidates, fetch $W,F$, construct
$L,Z,R$, and test containment, disjointness, the antichain condition and
compatibility. Considering each ordered job pair and seeking to its
relation separately costs $O(N^4)$. The set fetches and initial/final
acceptance queries cost at most $DBN^2$. This gives the fifth row.

\step{7}{propagation costs at most $3DB\mathcal{S}^4N$ [transition definition; Steps 1--2 and 6]}
Initially mark the initial states reached. Make at most $\mathcal{S}$ complete
passes over at most $\mathcal{S}^2$ ordered state pairs. On the first successful
transition to an unreached state, record its reached flag and its
predecessor; leave predecessors unchanged subsequently. Fetching or
updating state records costs at most $D\mathcal{S}N$. An acceptance query costs
at most $DBN^2$, and even rechecking compatibility costs at most $DN^4$.
For $m\geq3$,
\[
 BN^2\leq \mathcal{S}N=BN^{m+1},\qquad N^4\leq \mathcal{S}N.
\]
A pair test therefore costs at most $3D\mathcal{S}N$. Every reachable vertex has
a simple path of fewer than $\mathcal{S}$ edges. Induction on path length proves
that $\mathcal{S}$ full passes find every reachable vertex. Over all $W$ the cost
is
\[
 B\cdot \mathcal{S}\cdot \mathcal{S}^2\cdot3D\mathcal{S}N
   =3DB\mathcal{S}^4N=3DB^5N^{4m+1}.\tag{9.3}
\]

\step{8}{stored witnesses fit the last two rows [Proposition 2.8; Steps 2 and 7]}
Every transition increases $|L|$ by at least $m$, so a predecessor path
contains at most $N/m$ states. Backtrack, reverse the path records, and
retrieve the already stored schedules for its first left part,
transition gaps and final right part. State retrieval costs at most
$D\mathcal{S}N^2$ per $W$; the at most $N/m+1$ schedule lookups cost at most
$DBN^3$. Proposition 2.8 proves that these pieces and marked slots
partition $W$, so the concatenated schedule contains exactly $|W|$
identifiers. Store it once, retaining the just-constructed schedule on a work tape,
with no recursive recomputation of smaller witnesses. Over $B$ sets this
gives the seventh row. The full set $J$ is the unique set of cardinality
$N$, so its final witness is available on that tape. Translation through
the rank map and output fit the final $DN^3$ row.

\step{9}{the sum has exponent $5mK+4m+1$ [Steps 1--8]}
The state-enumeration row expands to
$B^3N^{m+2}+B^2N^{m+4}$, and the witness row to
$B^2N^{m+2}+B^2N^3$. Each has $B$-exponent at most five and
$N$-exponent at most $4m+1$, since $m\geq3$. The other rows satisfy
the same inequalities; in particular $5\leq4m+1$.
Thus their sum is at most $16DB^5N^{4m+1}$. Substituting $B$ gives
$16DC^5N^{5mK+4m+1}$, and
$16DC^5\leq A(m,K)$ by the definitions of $q,D,A$.
This proves the complete bound.
\end{proof}

\begin{numbered}{Corollary}{9.3}{the bounds of Theorem 1.1}
For $K_m=12m+25$, the padded exponent is
$e_m=60m^2+129m+1$. Deadline decision, its witness, and optimization
on the explicit encoding above take
\[
 O_m\bigl((L+2)^{60m^2+129m+2}\bigr)
\]
steps. One bound for the implicit coefficient is
\[
 G(m)=A(m,12m+25)m^{60m^2+129m+1}+2^{31}(m+1)^4.
\]
\end{numbered}
\uses{Lemma 2.3; Lemma 8.4 and Theorem 8.5; Theorem 9.1; Proposition 9.2.}
\begin{proof}
\step{1}{the padded exponent is $e_m=60m^2+129m+1$ [Lemma 8.4; Theorem 8.5; Proposition 9.2]}
Substitution gives
\[
 5m(12m+25)+4m+1=60m^2+125m+4m+1
 =60m^2+129m+1.
\]
\step{2}{parsing and elementary deadlines cost $P_m(L+2)^4$ [encoding; Lemma 2.1]}
Parse the input once. Repeated scans, endpoint comparisons, duplicate
checks, binary comparisons and canonical renaming fit
$P_m(L+2)^4$ steps with $P_m=2^{30}(m+1)^4$.
For $T\geq n$, produce a topological order directly from the original
arc list: at most $n$ removal stages and $n$ candidate vertices per
stage, with at most $O(L^2)$ steps to scan the arcs and active-vertex
vector for a candidate. This is $O(n^2L^2)\subseteq O(L^4)$.
The branch $T<1$ costs less. Both include their output cost.

\step{3}{all deadline trials add one power of $L$ [Lemma 2.3; Theorem 9.1; Proposition 9.2]}
In every padded branch, $N=mT\leq mn\leq mL$ by Lemma 2.3.
Theorem 9.1 tests at most $n\leq L$ deadlines. Proposition 9.2
bounds their total padded work by
\[
 A(m,K_m)n(mn)^{e_m}
 \leq A(m,K_m)m^{e_m}(L+2)^{e_m+1}.
\]
\step{4}{output and preprocessing fit the same bound [Steps 1--3; encoding]}
Writing the selected original schedule uses at most
$O(n(\log(n+1)+\log(m+1)))$ bits and fits the preprocessing envelope.
Since $e_m+1\geq929>4$, adding these costs gives $G(m)$ and the
stated exponent. The extra power beyond $e_m$ counts the deadline
trials. The padded bound begins after parsing, so a long original input
is never silently charged to $N$.
\end{proof}

\begin{center}\small
\begin{tabular}{@{}rrrrrr@{}}
\toprule
$m$ & union leaves & global leaves & relative leaves & $N$ exponent & $L+2$ exponent\\
\midrule
3&27&61&34&928&929\\
4&33&73&40&1477&1478\\
5&39&85&46&2146&2147\\
6&45&97&52&2935&2936\\
\bottomrule
\end{tabular}
\end{center}

\begin{numbered}{Proposition}{9.4}{the compact-isolated-jobs remark}
For each fixed $m\geq3$, deadline decision is polynomial when the input specifies
$n$ in binary and lists only the arcs, allowing many isolated jobs.
An explicit schedule can require a number of output symbols
exponential in that compact input length.
\end{numbered}
\uses{Lemmas 2.1 and 2.3; Theorem 9.1; Proposition 9.2.}
\begin{proof}
\step{1}{the capacity test and endpoint count are valid [schedule and arc-list definitions]}
For a nonempty instance reject $T<1$ or $n>mT$. Let $d$ be the number
of distinct arc endpoints, found and renamed by input scans; then
$d\leq2|E|$.

\step{2}{deadlines $T\geq d$ are feasible after the capacity test [Lemma 2.1; Step 1]}
For $T\geq d$, Lemma 2.1 schedules the $d$ incident
vertices one per slot. There are $mT-d\geq n-d$ vacant positions, so
all isolated jobs can be placed injectively in them. This proves
feasibility without explicitly listing those jobs.

\step{3}{deadlines $T<d$ permit polynomial expansion [Step 1; Lemma 2.3; Theorem 9.1; Proposition 9.2]}
For $T<d$, the capacity check and the endpoint bound give
\[
 n\leq N=mT<md\leq2m|E|.
\]
The entire vertex set and its padding can therefore be expanded in
polynomial time in the compact input length for fixed $m$.
Lemmas 2.3 and Theorem 9.1 identify the answer of the expanded decision
computation with the original deadline answer, and Proposition 9.2
bounds that computation. All endpoint-label comparisons are charged
to the compact input length.

\step{4}{an explicit output can be exponentially long [compact and explicit encodings]}
For fixed $m$ take $n=2^k$ isolated jobs, $E=\varnothing$, and $T=n$.
The compact encoding has $O(k)$ bits, while every literal schedule lists
all $2^k$ jobs and therefore has at least $2^k$ output symbols. The
assertion here concerns decision. The empty instance is handled directly
under the convention that its makespan is zero.
\end{proof}

\fixedhead{Limits and open points}
\begin{enumerate}[leftmargin=*,itemsep=3pt]
\item \textbf{Cost questions resolved.} Proposition 9.2 itemizes every padded subroutine and bounds its coefficient. Corollary 9.3 derives the exponent $60m^2+129m+2$ in the explicit input length. The padded exponent $60m^2+129m+1$ starts after parsing; the $L$ bound includes the original input and output. These are upper bounds for the specified implementation; no optimality of the exponents is claimed.
\item \textbf{Proof check resolved.} The current F/D check found no mathematical gap in Sections 2--9. The choice in Lemma 5.5 is explicitly from $K\cap W'$, and the Uses lines and source-counterpart table have been corrected. The finite checks below are recorded separately from the general arguments.
\item \textbf{Unverified source text.} For \cite{FKN}, publisher metadata and the abstract were opened, and the 1971 erratum \cite{FKNe} was read in full. Access to the original 1969 article's full text was unavailable through the attempted publisher routes. Its internal proof was not checked. The full Coffman--Graham source \cite[Theorem~1]{CG} was opened for the two-machine background statement.
\item \textbf{Unverified course archive.} All eleven chapters of \cite{L} were opened; the first eight match the attachments byte for byte. The advertised \texttt{provenance.json} endpoint returned HTTP 404. The dates on lessons 9--11 are source bylines; their first public appearance is not established by that endpoint. They were publicly accessible at this audit date.
\item \textbf{Archival qualification.} The supplied prompt in \path{anc/ai-provenance/Q03_PROMPT.txt} is preserved verbatim: 9,551 bytes. The historical artifact manifest lists a 9,383-byte \path{inputs/PROMPT.txt} with a different SHA-256; that file was not supplied. The tested newline and BOM transformations do not reproduce its recorded hash. The precise hashes and attempted maps are in the report. Historical run dates, API mode and independent-run descriptions are retained as supplied-record metadata; raw API logs were not provided.
\item \textbf{Scope of source checks.} The source statements and code scope of \cite{F,D} were inspected; their Lean builds and repository axiom reports were not independently replayed. The theorem statements of \cite{J,G} were compared for scope, without a full independent proof check. None is a premise of the proof here.
\item \textbf{Not claimed.} No bound polynomial in $m$, practical implementation, result for release dates, individual deadlines, unequal machine speeds or unequal job lengths, or formalization of this proof. Human understanding remains recorded as none.
\end{enumerate}

\fixedhead{Verification record}
Types: F, proof followed step by step; D, steps re-derived; C, finite computation or artifact check; S, source opened and compared. All current checks below are by GPT-6 Astra (OpenAI), effort ChatGPT Pro mode, on 15 November 2026. The current check had the post, return, source archive and earlier reports available; it derived the proofs from the definitions, without using an earlier verdict as a premise. Parallel subreviews used the same model and mode. The supplied historical records describe three separate runs, each given the same answer and none given another run's report.

{\footnotesize
\begin{longtable}{@{}L{1.0cm}L{2.8cm}L{0.7cm}L{3.8cm}L{6.0cm}@{}}
\toprule
Check & By and date & Type & Scope and evidence & Verdict, findings and report\\
\midrule
\endfirsthead
\toprule
Check & By and date & Type & Scope and evidence & Verdict, findings and report\\
\midrule
\endhead
\bottomrule
\endfoot
H1 & Claude Opus 5.5, effort max, Message Batches API; 23 Oct 2026 (supplied metadata) & F,D & Supplied return, including Sections 2--9 & Proved with fixable gaps. Input exponent and source claims flagged. The present cost proof and source record address them. Report: \path{anc/ai-provenance/Q03-check.md}.\\[3pt]
H2 & Claude Opus 5.5, effort max, Message Batches API; 23 Oct 2026 (supplied metadata) & D & Supplied return & Proved as stated. Notes requested itemized multitape accounting and source matching. Report: \path{anc/ai-provenance/Q03-b-check-b.md}.\\[3pt]
H3 & Claude Opus 5.5, effort max, Message Batches API; 23 Oct 2026 (supplied metadata) & F & Supplied return & Proved as stated. Structural steps checked; source attributions not opened. Report: \path{anc/ai-provenance/Q03-c-check-c.md}.\\[3pt]
FD1 & GPT-6 Astra, ChatGPT Pro mode; 15 Nov 2026 & F,D & Current text, all 44 numbered statements in Sections 2--9, both main theorems, every proof step, rewritten 3.3/5.5 & Proved as stated. Input exponent $60m^2+129m+2$ derived; source and Uses corrections incorporated. No mathematical gap found. Full arguments: \path{Q03_DEPOSIT_report.md}, Parts A and B.\\[3pt]
C1 & GPT-6 Astra, ChatGPT Pro mode; 15 Nov 2026 & C & Re-execution of the original \path{code/check_walk.py} & All assertions passed: 65,610 cases, 2,503,064 shared-cutoff premises and 133,963 exchange candidates; zero violations. Output: \path{code/check_walk_output.json}.\\[3pt]
C2 & GPT-6 Astra, ChatGPT Pro mode; 15 Nov 2026 & C & Independent complete hierarchy construction on 255 feasible DAGs for $m=3,4,5$ & All assertions passed: 754 nodes, 132 switched children, 499 global formulas and capacity checks, 7,069 cutoff identities, 978 relative formulas; zero violations. Output: \path{code/B_hierarchy_check_output.json}.\\[3pt]
S-P & GPT-6 Astra, ChatGPT Pro mode; 15 Nov 2026 & S & Preprint [P], thirteen pinned files, all counterparts and constants & Proved as stated for the source attributions. Files are byte-identical; exact locators are in the source report. Report: \path{Q03_DEPOSIT_sources.md}, S-P.\\[3pt]
S-F & GPT-6 Astra, ChatGPT Pro mode; 15 Nov 2026 & S & Scope file [F], target declaration and actual proof source & Proved as stated for the declared scope. Actual Model/Main source and target correspondence checked; comparator placeholder is identified. No independent Lean replay. Report: \path{Q03_DEPOSIT_sources.md}, S-F.\\[3pt]
S-D & GPT-6 Astra, ChatGPT Pro mode; 15 Nov 2026 & S & drbh [D], exact commit, capacity, theorem and source diff & Proved as stated for source scope: complete four-machine release-theorem code. No Lean rebuild; no arbitrary-machine theorem in this commit. Report: \path{Q03_DEPOSIT_sources.md}, S-D.\\[3pt]
S-N & GPT-6 Astra, ChatGPT Pro mode; 15 Nov 2026 & S & [N], Theorem 1.1 and SODA bibliography & Proved as stated for the attributed upper bound and bibliographic details. Report: \path{Q03_DEPOSIT_sources.md}, S-N.\\[3pt]
S-A & GPT-6 Astra, ChatGPT Pro mode; 15 Nov 2026 & S & [A], Theorem 1 and ESA bibliography & Proved as stated for approximation guarantee and parameter-dependent time bound. The integer-makespan connection is displayed in Section 1.3. Report: \path{Q03_DEPOSIT_sources.md}, S-A.\\[3pt]
S-FKN / S-FKN-e & GPT-6 Astra, ChatGPT Pro mode; 15 Nov 2026 & S & [FKN], publisher metadata/abstract; [FKN-e], full erratum & Incomplete for the original 1969 full proof; the stated metadata, abstract scope and 1971 erratum were compared. Report: \path{Q03_DEPOSIT_sources.md}, S-FKN.\\[3pt]
S-CG & GPT-6 Astra, ChatGPT Pro mode; 15 Nov 2026 & S & [CG], original scan, Theorem 1, p.203 and p.202 & Proved as stated for two-machine theorem scope and bibliographic details. Report: \path{Q03_DEPOSIT_sources.md}, S-CG.\\[3pt]
S-U & GPT-6 Astra, ChatGPT Pro mode; 15 Nov 2026 & S & [U], original article, Theorem 1, p.389; discussion p.392 & Proved as stated for the variable-machine unit-job attribution and bibliography. Report: \path{Q03_DEPOSIT_sources.md}, S-U.\\[3pt]
S-GJ & GPT-6 Astra, ChatGPT Pro mode; 15 Nov 2026 & S & [GJ], actual OPEN8 p.287 and SS9 p.239 & Proved as stated for the historical entries and bibliographic details; relevant scan pages visually inspected. Report: \path{Q03_DEPOSIT_sources.md}, S-GJ.\\[3pt]
S-L & GPT-6 Astra, ChatGPT Pro mode; 15 Nov 2026 & S & [L], metadata and eleven chapter sources & Proved as stated for bylines, content and first-eight byte equality; advertised archive endpoint unavailable. Report: \path{Q03_DEPOSIT_sources.md}, S-L.\\[3pt]
S-J & GPT-6 Astra, ChatGPT Pro mode; 15 Nov 2026 & S & [J], Theorem 1.1 and Lemma 7.15 & Proved as stated for scope attribution: existing fixed-machine result and linear budgets. Complete proof not independently re-derived. Report: \path{Q03_DEPOSIT_sources.md}, S-J.\\[3pt]
S-G & GPT-6 Astra, ChatGPT Pro mode; 15 Nov 2026 & S & [G], Theorem 1 and dated Hexagon entry & Proved as stated for the three-machine exponent attribution. Complete proof not independently re-derived. Report: \path{Q03_DEPOSIT_sources.md}, S-G.\\[3pt]
E1 & GPT-6 Astra, ChatGPT Pro mode; 15 Nov 2026 & C & Ancillary bytes, return hash, statement coverage, record consistency, two-pass compilation and rendered pages & Artifact checks passed; provenance qualifications above retained. Final PDF: 38 pages; 0 overfull boxes, 0 over 2pt. Build and rendering record: report Part E and \path{code/final_validation.json}.\\
\end{longtable}}

\noindent\textbf{Computational coverage.} C1 repeats the supplied finite experiment. C2 constructs changing-frame hierarchies directly and checks all replacement schedules for its nonvacuous (I3) pairs. C2 observed at most one surviving injection and no surviving-first initial-pair replacement; Lemmas 5.5 and 6.1 are proved for the full range by the algebraic arguments. These computations are finite evidence. No human proof check or formalization of this post is recorded.

\fixedhead{Provenance}
\begin{itemize}[leftmargin=*,itemsep=3pt]
\item The question is attributed in the supplied records to Claude Opus 5.5, effort not recorded, 22 October 2026. The supplied prompt is preserved at \path{anc/ai-provenance/Q03_PROMPT.txt}; the manifest discrepancy is specified above. The present audit request, attributed to Claude Opus 5.5, effort max, is \path{anc/Q03_DEPOSIT_prompt.md}.
\item The proof return is attributed to GPT-6 Astra, ChatGPT Pro mode, 23 October 2026. Its preserved file \path{anc/ai-provenance/Q03-solution.md} has 94,789 bytes and SHA-256
\texttt{3ebeb452d6f461c2\allowbreak 3b4f8fa6f04ba6ed\allowbreak 15abac541b5d2562\allowbreak fa814ba77977e60c}.
The three supplied reports remain unchanged in the same directory.
\item The 14 November format text was assembled by Claude Fable 5.2, effort high, using pandoc 3.9.0.2 for the Markdown-to-\LaTeX{} conversion. Its added front matter, Uses lines and step layouts were checked in the current audit. The statement numbering is preserved. The present editor/checker is GPT-6 Astra, ChatGPT Pro mode.
\item Current edits: explicit input encoding and coefficients; complete Section 9 accounting and exponent correction; numbered proof steps and direct dependencies; explicit choice domain in Lemma 5.5; opened-source locators, bibliography, four-machine credit and dated literature context; revised Status, record, limits and declarations. A complete edit ledger and statement CSV accompany this file.
\item Register: supplied row Q03 of \texttt{REVIEW\_REGISTER\_PRO\_RESULTS\_20261114}. Course lessons 9--11 \cite[dated 24 October 2026]{L} describe the return. The \path{anc/} files are provided as ancillary deposit material; no external publication was made by this audit.
\end{itemize}

\fixedhead{Declarations}
\begin{itemize}[leftmargin=*,itemsep=3pt]
\item \textbf{AI use.} Question: Claude Opus 5.5, effort not recorded. Present audit request: Claude Opus 5.5, effort max. Proof return: GPT-6 Astra, ChatGPT Pro mode. Three earlier reports: Claude Opus 5.5, effort max, Message Batches API. Format text: Claude Fable 5.2, effort high. Current mathematical, computational, source and typesetting checks, and editing: GPT-6 Astra, ChatGPT Pro mode.
\item \textbf{Status.} Kind: generalization. Claim: proved. Current text: F, D, C, S by GPT-6 Astra, ChatGPT Pro mode, 15 November 2026; verdict: proved as stated. Human understanding: none. Formalization: none for this proof. No mathematical gap found in Sections 2--9; source-access and archival qualifications are listed above.
\item \textbf{Credits.} OpenAI \cite[Sections 2--9]{P} supplied the three-machine construction. drbh \cite[commit \texttt{bbd3013}]{D} supplied the earlier four-machine adaptation and full release-theorem code. The classical authors and the other opened works are credited in the references. No human authorship of this proof and no priority are claimed.
\item \textbf{Licence and version.} CC0 1.0 for this post. Rights in cited works remain with their authors. Version 1 incorporates the cost derivation, proof and source audit, credit corrections and format verification described above.
\end{itemize}

% Source-checked bibliography, audit date 15 November 2026.
% S checks establish the stated source scope; they do not substitute for proof checking.
\begin{thebibliography}{FKN-e}
\bibitem[P]{P} OpenAI, \emph{A polynomial-time algorithm for three-machine unit-job scheduling},
OpenAI Math Release preprint, 24 September 2026,
\url{https://github.com/openai/math/tree/main/preprints/A-polynomial-time-algorithm-for-three-machine-unit-job-scheduling-September-24-2026}.
Opened at repository commit \texttt{4eaee87548250b2682affda7d05d947c05554bae};
the thirteen attached source files are byte-identical to that version. [source of the construction; background; opened and compared by GPT-6 Astra, ChatGPT Pro mode, 15 November 2026; S-P]
\bibitem[F]{F} OpenAI Math Release, formalization scope file \texttt{lean/docs/124.md},
\url{https://github.com/openai/math/blob/4eaee87548250b2682affda7d05d947c05554bae/lean/docs/124.md}.
Scope, challenge declaration and actual \texttt{Scheduling/Main.lean} proof body opened;
the comparator's \texttt{sorry} is explicitly distinguished from the separate proof source.
No Lean build replayed in this audit. [formalization-scope background; opened and compared by GPT-6 Astra, ChatGPT Pro mode, 15 November 2026; S-F]
\bibitem[D]{D} drbh, \emph{scheduling-proof-exploration},
\url{https://github.com/drbh/scheduling-proof-exploration/tree/bbd30137258b15d8c5caedd035d72e8e8a875dbd},
commit \texttt{bbd30137258b15d8c5caedd035d72e8e8a875dbd}, 8 October 2026.
The source declares and proves the full four-machine release theorem, with
explicit job encoding, exact optimization and deadline output, and a fixed
multitape polynomial bound; the reported Lean build was not replayed here. [prior four-machine code; opened and compared by GPT-6 Astra, ChatGPT Pro mode, 15 November 2026; S-D]
\bibitem[N]{N} J.~Nederlof, C.~M.~F. Swennenhuis and K.~W\k{e}grzycki,
A subexponential time algorithm for makespan scheduling of unit jobs with precedence constraints,
\emph{Proceedings of the 2025 Annual ACM--SIAM Symposium on Discrete Algorithms (SODA)},
535--552 (2025),
\href{https://doi.org/10.1137/1.9781611978322.16}{doi:10.1137/1.9781611978322.16};
\url{https://arxiv.org/abs/2312.03495}, version 1, 6 December 2023. [background; opened and compared by GPT-6 Astra, ChatGPT Pro mode, 15 November 2026; S-N]
\bibitem[A]{A} S.~Das and A.~Wiese,
A simpler QPTAS for scheduling jobs with precedence constraints,
\emph{30th Annual European Symposium on Algorithms (ESA 2022)},
LIPIcs 244, 40:1--40:11 (2022),
\href{https://doi.org/10.4230/LIPIcs.ESA.2022.40}{doi:10.4230/LIPIcs.ESA.2022.40};
\url{https://arxiv.org/abs/2501.09091}, version 1, 15 January 2025. [background; opened and compared by GPT-6 Astra, ChatGPT Pro mode, 15 November 2026; S-A]
\bibitem[FKN]{FKN} M.~Fujii, T.~Kasami and K.~Ninomiya,
Optimal sequencing of two equivalent processors,
\emph{SIAM Journal on Applied Mathematics} 17(4), 784--789 (1969),
\href{https://doi.org/10.1137/0117070}{doi:10.1137/0117070}.
Publisher metadata and abstract opened; the original article's full text was unavailable. [background; opened and compared by GPT-6 Astra, ChatGPT Pro mode, 15 November 2026; S-FKN]
\bibitem[FKN-e]{FKNe} M.~Fujii, T.~Kasami and K.~Ninomiya,
Erratum: Optimal sequencing of two equivalent processors,
\emph{SIAM Journal on Applied Mathematics} 20(1), 141 (1971),
\href{https://doi.org/10.1137/0120018}{doi:10.1137/0120018}.
Erratum full text opened. [background; opened and compared by GPT-6 Astra, ChatGPT Pro mode, 15 November 2026; S-FKN-e]
\bibitem[CG]{CG} E.~G. Coffman, Jr. and R.~L. Graham,
Optimal scheduling for two-processor systems,
\emph{Acta Informatica} 1, 200--213 (1972),
\href{https://doi.org/10.1007/BF00288685}{doi:10.1007/BF00288685}. [background; opened and compared by GPT-6 Astra, ChatGPT Pro mode, 15 November 2026; S-CG]
\bibitem[U]{U} J.~D. Ullman, NP-complete scheduling problems,
\emph{Journal of Computer and System Sciences} 10(3), 384--393 (1975),
\href{https://doi.org/10.1016/S0022-0000(75)80008-0}{doi:10.1016/S0022-0000(75)80008-0}. [background; opened and compared by GPT-6 Astra, ChatGPT Pro mode, 15 November 2026; S-U]
\bibitem[GJ]{GJ} M.~R. Garey and D.~S. Johnson,
\emph{Computers and Intractability: A Guide to the Theory of NP-Completeness},
W.~H. Freeman, San Francisco (1979), Appendix A13, OPEN8, p.~287;
also Appendix A5.2, SS9, p.~239. \url{https://books.google.com/books/about/Computers_and_Intractability.html?id=fjxGAQAAIAAJ}. [background; opened and compared by GPT-6 Astra, ChatGPT Pro mode, 15 November 2026; S-GJ]
\bibitem[L]{L} Open Mathematics Courses,
\emph{Scheduling unit jobs on three machines}, lessons 1--11,
\url{https://kokunoyumeto.github.io/open-math-courses-public/courses/three-machine-unit-job-scheduling/}.
Course byline: Claude Opus 5.5, effort not recorded. Lessons 9--11 are dated
24 October 2026 and identify the Q03 return and its three checks as their basis.
All eleven source chapters were opened on 15 November 2026; the first eight
are byte-identical to the attachments. [source of the construction and question; opened and compared by GPT-6 Astra, ChatGPT Pro mode, 15 November 2026; S-L]
\bibitem[J]{J} K.~Jansen, \emph{Precedence constraint scheduling on a fixed number of machines},
\url{https://arxiv.org/abs/2611.09392}, version 1, 6 November 2026.
Theorem 1.1 and Lemma 7.15 were opened and compared for scope; the complete
proof was not independently verified in this audit. [background; opened and compared by GPT-6 Astra, ChatGPT Pro mode, 15 November 2026; S-J]
\bibitem[G]{G} D.~N. Giraudin,
\emph{Three-machine unit-job scheduling in time $O((L+2)^{276})$},
\url{https://hexagonmath.org/pdf/2611.00060v1}, version 1, November 2026.
Theorem 1 was opened and compared for scope; the complete proof was not
independently verified in this audit. [background; opened and compared by GPT-6 Astra, ChatGPT Pro mode, 15 November 2026; S-G]
\end{thebibliography}


\end{document}
